\documentclass[11pt,a4paper]{article}

\usepackage[utf8]{inputenc}
\usepackage[T1]{fontenc}
\usepackage{amsmath,amssymb,amsthm}
\usepackage{bm}
\usepackage{mathtools}
\usepackage{geometry}
\geometry{margin=1in}
\usepackage{hyperref}
\usepackage{enumitem}

\DeclareMathOperator{\SO}{SO}
\DeclareMathOperator{\softmax}{softmax}
\DeclareMathOperator{\softplus}{softplus}
\DeclareMathOperator{\sinc}{sinc}
\DeclareMathOperator{\diag}{diag}
\DeclareMathOperator{\Span}{span}
\DeclareMathOperator{\sgn}{sgn}
\DeclareMathOperator{\ScalarRMSNorm}{ScalarRMSNorm}
\DeclareMathOperator{\SiLU}{SiLU}
\DeclareMathOperator*{\SUM}{\textstyle\sum}

\newcommand{\R}{\mathbb{R}}
\newcommand{\N}{\mathbb{N}}
\newcommand{\Z}{\mathbb{Z}}
\newcommand{\C}{\mathbb{C}}
\newcommand{\bA}{\bm{A}}
\newcommand{\bU}{\bm{U}}
\newcommand{\bV}{\bm{V}}
\newcommand{\bG}{\bm{G}}
\newcommand{\Sph}{S^2}
\newcommand{\bR}{\bm{R}}
\newcommand{\br}{\bm{r}}
\newcommand{\brh}{\hat{\bm{r}}}
\newcommand{\bx}{\bm{x}}
\newcommand{\by}{\bm{y}}
\newcommand{\bz}{\bm{z}}
\newcommand{\bmF}{\bm{F}}
\newcommand{\bh}{\bm{h}}
\newcommand{\bq}{\bm{q}}
\newcommand{\bk}{\bm{k}}
\newcommand{\bv}{\bm{v}}
\newcommand{\bg}{\bm{g}}
\newcommand{\bu}{\bm{u}}
\newcommand{\Te}{\bm{T}}
\newcommand{\bD}{\bm{D}}
\newcommand{\bW}{\bm{W}}
\newcommand{\bP}{\bm{P}}
\newcommand{\bphi}{\bm{\phi}}
\newcommand{\brho}{\bm{\rho}}
\newcommand{\rcut}{r_{\mathrm{c}}}
\newcommand{\rin}{r_{\mathrm{in}}}
\newcommand{\rout}{r_{\mathrm{out}}}
\newcommand{\lmax}{l_{\max}}
\newcommand{\mmax}{m_{\max}}
\newcommand{\rtilde}{\widetilde{r}}
\newcommand{\Nb}{\mathcal{N}}
\newcommand{\Idx}{\mathcal{I}}

\theoremstyle{plain}
\newtheorem{theorem}{Theorem}
\newtheorem{proposition}[theorem]{Proposition}
\newtheorem{lemma}[theorem]{Lemma}
\theoremstyle{definition}
\newtheorem{definition}[theorem]{Definition}
\theoremstyle{remark}
\newtheorem{remark}[theorem]{Remark}

\title{The DPA4 Architecture:\\ A Mathematical Construction}
\author{}
\date{}

\begin{document}
\maketitle

\begin{abstract}
    We construct the DPA4 message-passing equivariant graph neural network as a mathematical object: a parametric family of $C^3$ scalar functions of atomic configurations that is invariant under translations, species-preserving permutations, and global rotations, and that admits a strict freezing property in a designated short-range bridging zone.
    The construction proceeds by introducing, in order, the configuration space, the symmetry representation, the per-edge local frame, the bridging window, the equivariant feature space, the equivariant operators, the interaction block, and the energy readout.
    Every symbol is defined at its first occurrence.
\end{abstract}

\tableofcontents

% =============================================================
\section{Overview of the Architecture}\label{sec:overview}

DPA4 is a parametric family of $C^3$ scalar functions
\begin{equation}\label{eq:overall}
    E_\Theta : (Z,R)\;\longmapsto\;E_\Theta(Z,R)\in\R
\end{equation}
of atomic configurations $(Z,R)$ (species $Z\in\{1,\dots,n_t\}^N$, positions $R\in\R^{3N}$), invariant under translations, species-preserving permutations, and global rotations, and conservative in the sense that forces and stress are exact gradients of $E_\Theta$.
We now sketch the construction; the formal definitions appear in the following sections.

\paragraph{Two additive branches.}
The energy decomposes as
\begin{equation}\label{eq:overview-decomp}
    E_\Theta(Z,R) \;=\; E^{\mathrm{NN}}_\Theta(Z,R) \;+\; E^{\mathrm{ZBL}}(Z,R),
\end{equation}
where $E^{\mathrm{ZBL}}$ is a fixed analytical short-range pair potential (\S\ref{sec:zbl}) and $E^{\mathrm{NN}}_\Theta$ is the learned message-passing graph neural network branch defined by the remainder of the document.
The geometric inputs consumed by the two branches differ inside a designated short-range \emph{bridging window} $[\rin,\rout]$: $E^{\mathrm{NN}}_\Theta$ reads a smoothly clamped distance $\rtilde(r)\equiv\rin$ for $r\leq\rin$, while $E^{\mathrm{ZBL}}$ reads the true $r$.
A multiplicative \emph{source-freeze gate} $\eta_j\in[0,1]$ further enforces that any atom $j$ which sees a partner inside the frozen zone contributes nothing to the learned branch (Theorem~\ref{thm:freeze}).

\paragraph{State space.}
The learned branch operates on per-atom states living in a finite-dimensional
representation of $\SO(3)$:
\begin{equation}\label{eq:overview-state}
    X\;=\;(\bh_1,\dots,\bh_N),\qquad
    \bh_i \;\in\; V_{\le\lmax}\otimes\R^{C},
    \quad
    V_{\le\lmax}\;=\;\bigoplus_{l=0}^{\lmax}V_l,
\end{equation}
where $V_l$ is the real $(2l+1)$-dimensional irreducible representation of $\SO(3)$ and $C\in\N$ is a channel width.
The $\SO(3)$ action $X\mapsto\bD(R)X$ is channel-wise.
The $l=0$ subspace is the trivial representation; energies will be read out from it alone.

\paragraph{From global SO(3) to per-edge SO(2).}
For each directed edge $(i,j)$ with $r_{ij}>0$ we construct a $C^\infty$
rotation $R_{ij}\in\SO(3)$ that aligns the edge direction with a fixed
reference vector $\bz=(0,0,1)^\top$, and let $\bD_{ij}\coloneqq\bD(R_{ij})$ be
the corresponding representation. Inside the edge-local frame the residual
symmetry is no longer $\SO(3)$ but its abelian subgroup $\SO(2)_{\bz}$
stabilising $\bz$. Under $\SO(2)_{\bz}$ the basis labelled by $(l,m)$
decomposes into one-dimensional invariant lines ($m=0$) and real
two-dimensional components ($|m|>0$). Each interaction therefore needs only a
block-diagonal $\SO(2)$-equivariant mixing along an angular order $m$, instead
of the full Clebsch--Gordan tensor product required by SO(3)-equivariant
operators. Restricting further to $|m|\leq\mmax$ for some $\mmax\leq\lmax$
yields a reduced layout of dimension $D_m=\sum_l(2\min(l,\mmax)+1)$, with a
per-degree rescale $\rho_l$ that compensates the truncation.

\paragraph{Interaction block.}
A single interaction block $\mathcal{B}$ maps a state to a state by two
equivariant residual updates. Writing $\bh_i'$ for the intermediate state after
the first update and $\bh_i''$ for the output of the block,
\begin{equation}\label{eq:overview-block}
    \bh_i' = \bh_i + \mathcal{C}(\bh)_i,\qquad
    \bh_i'' = \bh_i' + \mathcal{F}\bigl(\bh'\bigr)_i,
\end{equation}
where (i) $\mathcal{C}$ is the \emph{edge convolution} (\S\ref{sec:so2conv}): it pre-mixes channels, transports the source feature into the edge-local frame via $\bD_{ij}^{\le\mmax}$, modulates the result by per-edge per-degree radial features $\brho_{ij,l}$, applies a stack of SO(2)-equivariant linear/normalise/gated nonlinearity layers (optionally with cross-focus competition), lifts back to the global frame with the rescale $\rho_l$, and aggregates over neighbours through either an envelope-weighted sum or an envelope-gated softmax attention with a query-dependent output gate; (ii) $\mathcal{F}$ is an equivariant feed-forward block built from degree-wise linear maps and a gated activation that mixes $l=0$ scalars with higher-$l$ vectors through scalar gates.
Both $\mathcal{C}$ and $\mathcal{F}$ start as the identity at initialisation (zero-initialised output projections), so $\mathcal{B}$ is a near-identity perturbation of the input state at training time zero.
The primes in \eqref{eq:overview-block} denote \emph{intra-block} substates; \emph{inter-block} iterates are tracked separately by the counter $X^{(b)}$ introduced below.

\paragraph{Computational pipeline.}
For a given input $(Z,R)$, $E^{\mathrm{NN}}_\Theta$ is computed by the
following composition of maps:
\begin{equation}\label{eq:overview-pipeline}
    (Z,R)\;\xrightarrow{\text{edges}}\;
    \{(i,j,\widetilde\br_{ij},R_{ij},\eta_j,\dots)\}\;
    \xrightarrow{\text{init}}\;X^{(0)}\;
    \xrightarrow{\mathcal{B}_1}\cdots\xrightarrow{\mathcal{B}_B}\;X^{(B)}\;
    \xrightarrow{\mathcal{O}}\;\sum_{i=1}^{N}E_i^{\mathrm{NN}}.
\end{equation}
Here the superscript $(b)$ counts the \emph{number of interaction blocks already applied}:
\begin{equation}\label{eq:Xb-recursion}
    X^{(b)} \;=\; \mathcal{B}_b\bigl(X^{(b-1)}\bigr),\qquad b=1,\dots,B,
\end{equation}
so that $X^{(B)}$ is the final state read out by $\mathcal{O}$, and the \emph{initial state}
\begin{equation}
    X^{(0)} \;=\; \bigl(\bh_1^{(0)},\dots,\bh_N^{(0)}\bigr) \;\in\;\bigl(V_{\le\lmax}\otimes\R^C\bigr)^{N}
\end{equation}
is the input to the first block, produced by the initialisation step before any message passing.
The pieces of the pipeline are:
\begin{itemize}[leftmargin=2em]
    \item the \emph{edge construction} step (\S\ref{sec:radial},\S\ref{sec:sfpg})
          computes the clamped displacement $\widetilde\br_{ij}$, the edge-local rotation
          $R_{ij}$, the smooth cutoff factor $s(r_{ij})$, the source-freeze gate
          $\eta_j$, and a smooth in-degree $n_i$;
    \item the \emph{initialisation} step (\S\ref{sec:init}) defines $X^{(0)}$ explicitly:
          \begin{equation}\label{eq:overview-X0}
              \bh_i^{(0)}\bigm|_{l=0,m=0,c}\;=\;T_{Z_i,c},
              \qquad
              \bh_i^{(0)}\bigm|_{(l,m'),c}\;=\;n_i\!\!\sum_{j:r_{ij}<\rcut}\!\!\eta_j\,Y_l^{m'}(\brh_{ij})\,\brho_{ij,l,c}\ \ (l\geq 1),
          \end{equation}
          where $T\in\R^{n_t\times C}$ is a learnable species table, $Y_l^{m'}$ is the real spherical harmonic of degree $l$ and order $m'$ (the geometric initial embedding), and $\brho_{ij,l}$ are the per-edge per-degree radial features (\eqref{eq:rho-edge}); an optional FiLM modulation derived from a local environment matrix may additionally modulate the $l=0$ slice (\S\ref{sec:init});
    \item $B$ \emph{interaction blocks} $\mathcal{B}_b$ are applied in sequence via \eqref{eq:Xb-recursion}, with an optional non-increasing pyramid schedule $(\lmax^{(b)},\mmax^{(b)})$ that truncates the working representation between blocks (\S\ref{sec:energy});
    \item the \emph{scalar readout} $\mathcal{O}:\R^{C}\to\R$ is applied to the $l=0$
          slice of each final node feature $\bh_i^{(B)}$, yielding per-atom energies
          $E_i^{\mathrm{NN}}=\mathcal{O}(\bh_i^{(B)}|_{l=0})$ whose sum is
          $E^{\mathrm{NN}}_\Theta$.
\end{itemize}

\paragraph{Guarantees.}
By construction (see \S\ref{sec:thms}):
\begin{itemize}[leftmargin=2em]
    \item Every block preserves $\SO(3)$ equivariance of the state, and the readout is
          $\SO(3)$-invariant; hence $E^{\mathrm{NN}}_\Theta$ is invariant under
          translations, permutations, and rotations.
    \item Each elementary ingredient (envelope, radial, distance clamp, bridging switch,
          two-chart edge quaternion, softmax denominator) is $C^3$ on its domain; hence
          $E_\Theta\in C^3$ off the configuration-diagonal.
    \item For any pair $(j,k)$ with $r_{jk}\leq\rin$, the gate $\eta_j$ and $\eta_k$
          vanish exactly; combined with the distance clamp this forces $\partial
              E^{\mathrm{NN}}_\Theta/\partial\bR_j=\partial
              E^{\mathrm{NN}}_\Theta/\partial\bR_k=0$, so the force on each member of a
          frozen pair is the pure ZBL force.
\end{itemize}

The rest of the document develops each piece of \eqref{eq:overview-pipeline} as
a precise mathematical object.

% =============================================================
\section{Configuration Space and Symmetries}\label{sec:setup}

\paragraph{Configurations.}
A configuration of $N\in\N$ atoms is a pair $\bigl(Z,R\bigr)$ where
$Z=(Z_1,\dots,Z_N)\in\{1,\dots,n_t\}^N$ assigns to each atom a species label
drawn from a fixed alphabet of size $n_t\in\N$, and
$R=(\bR_1,\dots,\bR_N)\in\R^{3N}$ assigns Cartesian positions. Optional
periodic boundary conditions, encoded by a cell matrix $\bm{L}\in\R^{3\times
        3}$, are treated by interpreting the index $j$ in pairs $(i,j)$ as ranging over
images of the unit cell.

For ordered pairs $(i,j)$ with $i\neq j$ define
\begin{equation}
    \br_{ij}\coloneqq \bR_j-\bR_i,\qquad
    r_{ij}\coloneqq \|\br_{ij}\|,\qquad
    \brh_{ij}\coloneqq \br_{ij}/r_{ij}\ \text{when }r_{ij}>0.
\end{equation}

\paragraph{Three symmetry groups.}
Let $\R^3$ act by translations on positions, $\SO(3)$ act by global rotations,
and the symmetric group $S_N$ act by species-preserving permutations of atom
indices. A function $E:\R^{3N}\to\R$ on configurations is \emph{admissible} if
it is invariant under each of these actions and is $C^3$ on the open set where
$\bR_i\neq\bR_j$.

\paragraph{Total energy.}
The DPA4 energy decomposes as
\begin{equation}\label{eq:totalE}
    E\bigl(Z,R\bigr) \;=\; E^{\mathrm{NN}}\bigl(Z,R\bigr) \;+\; E^{\mathrm{ZBL}}\bigl(Z,R\bigr),
\end{equation}
where $E^{\mathrm{ZBL}}$ is the analytical Ziegler--Biersack--Littmark short-range branch defined in \S\ref{sec:zbl} and $E^{\mathrm{NN}}$ is the learned graph neural network branch constructed in the remainder of the paper.
Atomic forces and virial stress are obtained by automatic differentiation of $E$ with respect to the positions and the cell;
in particular, conservativeness is by construction.

% =============================================================
\section{Representation Theory of $\SO(3)$ in a Real Basis}\label{sec:rep}

\paragraph{Real irreducible representations.}
For each non-negative integer $l\in\Z_{\geq 0}$, let $V_l$ be a real
$(2l+1)$-dimensional vector space carrying the real form of the spin-$l$
irreducible representation of $\SO(3)$. We index a basis of $V_l$ by an integer
$m\in\{-l,\dots,+l\}$. The representation is the group homomorphism
\begin{equation}
    D^l : \SO(3)\;\longrightarrow\;\mathrm{O}\bigl(2l+1\bigr),\qquad D^l(R)\,D^l(R')=D^l(RR').
\end{equation}
Each $D^l$ is orthogonal: $D^l(R)^\top D^l(R)=I_{2l+1}$.

\paragraph{Truncated direct sum.}
Fix a degree cutoff $\lmax\in\N_0$ and set $D\coloneqq(\lmax+1)^2$. The total
feature space carries the representation
\begin{equation}\label{eq:Drep}
    V_{\le\lmax} \;\coloneqq\; \bigoplus_{l=0}^{\lmax}V_l \;\cong\;\R^{D},\qquad
    \bD(R) \;\coloneqq\; \bigoplus_{l=0}^{\lmax}D^l(R)\;\in\;\R^{D\times D}.
\end{equation}
We linearly pack a basis of $V_{\le\lmax}$ by $\iota(l,m)\coloneqq l^2+l+m$.
The matrix $\bD(R)$ is block-diagonal in this packing.

\paragraph{Equivariance.}
A map $\Phi:V_{\le\lmax}^{\otimes p}\to V_{\le\lmax}^{\otimes q}$ between
tensor powers is \emph{$\SO(3)$-equivariant} if
$\Phi(\bD(R)x_1,\dots,\bD(R)x_p)=\bD(R)\Phi(x_1,\dots,x_p)$ for all
$R\in\SO(3)$. It is \emph{$\SO(3)$-invariant} if its image lies in the $l=0$
subspace.

\paragraph{Feature space.}
Let $C\in\N$ be a channel width. A \emph{node feature} at atom $i$ is an
element $\bh_i\in V_{\le\lmax}\otimes\R^C\cong\R^{D\times C}$; $\SO(3)$ acts on
the first factor only. The $l$-block of $\bh_i$ is $\bh_i^{(l)}\in
    V_l\otimes\R^C$.

% =============================================================
\section{Local Frame on a Directed Edge}\label{sec:frame}

\paragraph{Goal.}
For each directed pair $(i,j)$ with $r_{ij}>0$ we choose a rotation that maps
$\brh_{ij}$ to a fixed reference vector $\bz\coloneqq(0,0,1)^\top\in\Sph$. This
is exactly the data of a section $\sigma:\Sph\setminus\{*\}\to\SO(3)$ of the
bundle $\SO(3)\to\Sph$ given by the orbit map; the obstruction to a global
section is one point on $\Sph$ (the Euler-angle singularity). DPA4 uses two
charts and blends them with a $C^\infty$ partition of unity to obtain a
globally smooth lift.

\paragraph{Quaternion charts.}
Let $S^3\coloneqq\{\bq\in\R^4:\|\bq\|=1\}$ be the unit $3$-sphere of $\R^4$,
identified with the group of unit quaternions
$\bq=q_w+q_x\,\mathbf{i}+q_y\,\mathbf{j}+q_z\,\mathbf{k}$ in the components
$(q_w,q_x,q_y,q_z)\in\R^4$. The standard double cover $R:S^3\to\SO(3)$, under
which $\bq$ and $-\bq$ represent the same rotation, gives the identification
$\SO(3)\cong S^3/\{\pm 1\}$. Explicitly, $R$ is the map
\begin{equation}\label{eq:R-quat}
    R(\bq) \;=\;
    \begin{pmatrix}
        1-2(q_y^2+q_z^2) & 2(q_xq_y-q_wq_z) & 2(q_xq_z+q_wq_y) \\
        2(q_xq_y+q_wq_z) & 1-2(q_x^2+q_z^2) & 2(q_yq_z-q_wq_x) \\
        2(q_xq_z-q_wq_y) & 2(q_yq_z+q_wq_x) & 1-2(q_x^2+q_y^2)
    \end{pmatrix}\;\in\;\SO(3),
\end{equation}
or, equivalently, the conjugation action $R(\bq)\bm{v}=\bq\,\bm{v}\,\bq^{*}$ on the embedding $\R^3\hookrightarrow\mathbb{H}$, $\bm{v}\mapsto v_x\mathbf{i}+v_y\mathbf{j}+v_z\mathbf{k}$, where $\bq^{*}=q_w-q_x\mathbf{i}-q_y\mathbf{j}-q_z\mathbf{k}$ is the quaternion conjugate.

For $\brh=(x,y,z)\in\Sph$ define
\begin{equation}\label{eq:charts}
    \bq^{+}(\brh)\;=\;\frac{(1+z,\;y,\;-x,\;0)}{\sqrt{2(1+z)}},
    \qquad
    \bq^{-}(\brh)\;=\;\frac{(-x,\;0,\;1-z,\;y)}{\sqrt{2(1-z)}},
\end{equation}
defined respectively for $z>-1$ and $z<+1$.
Each is a smooth lift on its domain satisfying $R(\bq^{\pm}(\brh))\,\brh=\bz$, where $R:S^3\to\SO(3)$ is the standard double cover.

\paragraph{Blending.}
Let $\lambda(\brh)\coloneqq(1+z)/2\in[0,1]$ and choose the sign of $\bq^{-}$ to
make $\langle\bq^{+},\bq^{-}\rangle\geq 0$ (shortest arc). Define
\begin{equation}\label{eq:nlerp}
    \bq_{ij} \;\coloneqq\; \frac{\lambda\,\bq^{+}(\brh_{ij}) + (1-\lambda)\,\bq^{-}(\brh_{ij})}{\bigl\|\lambda\,\bq^{+}(\brh_{ij}) + (1-\lambda)\,\bq^{-}(\brh_{ij})\bigr\|}.
\end{equation}
The denominator is bounded away from zero on the overlap by the shortest-arc choice, so $\brh\mapsto\bq_{ij}\in S^3$ is $C^\infty$ on all of $\Sph$.
Pulling $\bq_{ij}$ through $R$ yields a $C^\infty$ rotation
\begin{equation}
    R_{ij}\coloneqq R(\bq_{ij})\in\SO(3),\qquad R_{ij}\,\brh_{ij}=\bz.
\end{equation}

\paragraph{Edge-local action.}
The local representation acts on features by $\bD_{ij}\coloneqq\bD(R_{ij})$. By
construction $\bD_{ij}$ fixes the direction $\bz$ pointwise; the residual
symmetry stabilising $\bz$ is the subgroup $\SO(2)_{\bz}\subset\SO(3)$
generated by rotations about $\bz$. Working in the real basis
$\{\mathbf{e}_m^{(l)}\}_{m=-l}^{+l}$ of $V_l$, the restriction of $V_l$ to
$\SO(2)_{\bz}$ decomposes into real irreducible subrepresentations as
\begin{equation}\label{eq:so2decomp}
    V_l\bigm|_{\SO(2)_{\bz}}\;=\;
    \underbrace{\R\,\mathbf{e}_0^{(l)}}_{\text{trivial $1$-d}}
    \;\oplus\;\bigoplus_{m=1}^{l}\mathcal{V}_m^{(l)},
    \qquad
    \mathcal{V}_m^{(l)}\;\coloneqq\;\R\,\mathbf{e}_{-m}^{(l)}\oplus\R\,\mathbf{e}_{+m}^{(l)},
\end{equation}
where $R_\theta\in\SO(2)_{\bz}$ acts trivially on $\R\,\mathbf{e}_0^{(l)}$ and on $\mathcal{V}_m^{(l)}$ ($m\geq 1$) in the basis $(\mathbf{e}_{-m}^{(l)},\mathbf{e}_{+m}^{(l)})$ by the planar rotation
\begin{equation}\label{eq:so2-rot}
    R_\theta\bigm|_{\mathcal{V}_m^{(l)}}\;=\;\begin{pmatrix}\cos m\theta & -\sin m\theta\\ \sin m\theta & \phantom{-}\cos m\theta\end{pmatrix}.
\end{equation}
Each $\mathcal{V}_m^{(l)}$ is a $2$-dimensional real irreducible representation of $\SO(2)_{\bz}$, isomorphic as a real $\SO(2)_{\bz}$-representation to the complex line $\C_m$ on which $R_\theta$ acts by multiplication by $e^{im\theta}$.
The total dimension is $1+2l=2l+1=\dim V_l$.
Individual real lines $\R\,\mathbf{e}_{\pm m}^{(l)}$ with $m\neq 0$ are \emph{not} invariant under $\SO(2)_{\bz}$ on their own; only the pair $\mathcal{V}_m^{(l)}$ is.

\paragraph{Remark: why no Clebsch--Gordan is needed.}
Clebsch--Gordan coefficients are the explicit basis of the intertwiner space
for a \emph{bilinear} $G$-equivariant product of irreps,
\[
    V_{a}\otimes V_{b}\longrightarrow V_{c},
\]
and they are non-trivial only when both factors are at $l>0$. For an
SO(3)-equivariant such product, the selection rule is $|a-b|\le c\le a+b$ and
the basis $\langle a\,m_a;b\,m_b\mid c\,m_c\rangle$ is a non-trivial table;
evaluating it at all admissible triples up to $\lmax$ costs
$\mathcal{O}(\lmax^6)$. This is the operation around which networks built on
full SO(3) irreps (e.g.\ NequIP, MACE) organise their cost: they assign to each
edge an \emph{equivariant} feature $Y_l^{m}(\brh_{ij})$ at $l=0,1,\dots,\lmax$
and then combine it with the source node feature $\bh_j$ via an SO(3) tensor
product, which forces a CG expansion.

DPA4 is structurally different: the only edge-side object ever multiplied into
an equivariant node feature is the \emph{$\SO(3)$-invariant} ($l=0$) radial
profile $\brho_{ij,l}$. The angular content of the edge geometry enters the
architecture not as an equivariant feature but as the rotation operator
$\bD_{ij}\in\bD(\SO(3))$ from \S\ref{sec:frame}, which acts on a single
equivariant feature by frame change ($\bh_j\mapsto\bD_{ij}\bh_j$), not as a
factor in a tensor product. Concretely:
\begin{itemize}[leftmargin=2em]
    \item In the geometric initial embedding (\S\ref{sec:init}, equation~\eqref{eq:gie}),
          the message is the product $Y_l^{m'}(\brh_{ij})\cdot\brho_{ij,l,c}$: an
          equivariant degree-$l$ vector multiplied by an invariant scalar — the trivial
          intertwiner $V_l\otimes\R\to V_l$.
    \item In the edge convolution (\S\ref{sec:so2conv}), the source feature is first
          rotated to the local frame, modulated per degree by the invariant scalar
          $\rho_l$, mixed by the SO(2)-equivariant linear operator $L^{\mathrm{SO2}}$,
          and rotated back: at no point are two equivariant features multiplied together.
    \item The gated nonlinearity $\Gamma_{\psi,\bG}$ multiplies higher-$l$ blocks by a
          sigmoid gate computed from the $\SO(3)$-invariant $l=0$ slice — again
          $V_l\otimes\R\to V_l$.
\end{itemize}
Every multiplicative combination is therefore of the trivial form "scalar $\times$ equivariant", whose intertwiner is one-dimensional and $G$-independent; no CG coefficient is required, in either $\SO(3)$ or $\SO(2)$.

The local-frame construction of \S\ref{sec:frame} plays a complementary role:
it makes the residual symmetry abelian, so even if one wanted to introduce a
bilinear equivariant product in the local frame, the Clebsch--Gordan series of
the abelian group $\SO(2)_{\bz}$,
\begin{equation}\label{eq:so2-cg}
    \C_{m_1}\otimes\C_{m_2}\;\cong\;\C_{m_1+m_2},
\end{equation}
degenerates to the Kronecker delta $\delta_{m_3,m_1+m_2}$ and no tabulated coefficient is needed.
The two facts together — equivariant edge features kept at $l=0$, and the abelian residual symmetry — yield a per-edge angular cost of $\mathcal{O}(\lmax\,\mmax\,C)$, in contrast to the $\mathcal{O}(\lmax^6)$ scaling of SO(3)-CG architectures.

% =============================================================
\section{Smooth Cutoff, Radial Basis, and Bridging Window}\label{sec:radial}

Let $\rcut>0$ be a fixed cutoff radius.

\paragraph{Cutoff envelope.}
For an integer exponent $p\in\N$ let $a_p,b_p,c_p,d_p\in\R$ be the unique
solution of $s_p(\rcut)=s_p'(\rcut)=s_p''(\rcut)=s_p'''(\rcut)=0$ for
\begin{equation}\label{eq:env}
    s_p(r) \coloneqq\begin{cases}
        1+x^p\bigl(a_p+b_p x+c_p x^2+d_p x^3\bigr), & x=r/\rcut\in[0,1), \\[2pt]
        0,                                          & x\geq 1.
    \end{cases}
\end{equation}
Solving yields $(a_p,b_p,c_p,d_p)=\bigl(-\tbinom{p+3}{3},3\tbinom{p+3}{2}/{(p+1)},\dots\bigr)$;
in closed form
\(
a_p=-\tfrac{(p+1)(p+2)(p+3)}{6}
\), \(
b_p=\tfrac{p(p+2)(p+3)}{2}
\), \(
c_p=-\tfrac{p(p+1)(p+3)}{2}
\), \(
d_p=\tfrac{p(p+1)(p+2)}{6}
\).
For $p=5$, $s_5(r)=1-56x^5+140x^6-120x^7+35x^8$;
for $p=7$, $s_7(r)=1-120x^7+540x^8-1080x^9+840x^{10}$.

\begin{lemma}\label{lem:envelope}
    $s_p\in C^3(\R)$, $s_p(0)=1$, $s_p(r)=0$ for $r\geq\rcut$, and all derivatives of $s_p$ of orders $0,1,2,3$ vanish at $r=\rcut$.
\end{lemma}

We henceforth fix two instances: $s\coloneqq s_5$ for per-edge message
weighting, and $\sigma_r\coloneqq s_7$ folded into the radial basis below.

\paragraph{Radial basis.}
Fix $n_r\in\N$ and a vector of positive frequencies
$\omega=(\omega_1,\dots,\omega_{n_r})\in\R_{>0}^{n_r}$. The radial basis is the
function $\bphi:\R_{\geq 0}\to\R^{n_r}$ with components
\begin{equation}\label{eq:rbf}
    \phi_n(r)\;\coloneqq\;\frac{\sin(\omega_n r)}{r}\cdot \sigma_r(r),\qquad n=1,\dots,n_r,
\end{equation}
with $\phi_n(0)=\omega_n\,\sigma_r(0)$ by continuous extension.
Each $\phi_n$ is real-analytic on $\R_{>0}$ and lies in $C^3(\R)$.
The frequencies $\omega$ are learnable; their initial values are $\omega_n^{(0)}=n\pi/\rcut$.

\paragraph{Bridging window.}
Choose $0<\rin<\rout$. Let $h_{\mathrm{c}},h_{\mathrm{w}}:[0,1]\to[0,1]$ be the
septic Hermite polynomials
\begin{align}
    h_{\mathrm{c}}(t) & \coloneqq 20t^4-45t^5+36t^6-10t^7,\label{eq:hclamp} \\
    h_{\mathrm{w}}(t) & \coloneqq 35t^4-84t^5+70t^6-20t^7.\label{eq:hsw}
\end{align}
The boundary conditions are tailored to the two distinct gluing requirements they will fulfil.
The clamp polynomial $h_{\mathrm{c}}$ must glue $\rin$ (constant) on the left to $r$ (identity) on the right, so its slope must equal the slopes of those two functions at the endpoints:
\begin{equation}\label{eq:hc-bcs}
    h_{\mathrm{c}}(0)=0,\ \ h_{\mathrm{c}}(1)=1,\ \ h_{\mathrm{c}}'(0)=0,\ \ h_{\mathrm{c}}'(1)=1,\ \ h_{\mathrm{c}}''(0)=h_{\mathrm{c}}''(1)=h_{\mathrm{c}}'''(0)=h_{\mathrm{c}}'''(1)=0.
\end{equation}
The switch polynomial $h_{\mathrm{w}}$ must glue the constants $0$ on the left and $1$ on the right, so \emph{all} of its first three derivatives vanish at both endpoints:
\begin{equation}\label{eq:hw-bcs}
    h_{\mathrm{w}}(0)=0,\ \ h_{\mathrm{w}}(1)=1,\ \ h_{\mathrm{w}}^{(k)}(0)=h_{\mathrm{w}}^{(k)}(1)=0\quad(k=1,2,3).
\end{equation}
These $4+4=8$ conditions on each side determine the unique septic interpolants \eqref{eq:hclamp}--\eqref{eq:hsw}.

Define the \emph{clamped distance map} $\rtilde:\R_{\geq 0}\to\R_{\geq 0}$ and
the \emph{bridging amplitude} $w:\R_{\geq 0}\to[0,1]$ by
\begin{equation}\label{eq:rtilde}
    \rtilde(r)\coloneqq\begin{cases}\rin,&r\leq\rin,\\ \rin+(\rout-\rin)\,h_{\mathrm{c}}\!\bigl(t(r)\bigr),&\rin<r<\rout,\\ r,&r\geq\rout,\end{cases}
    \qquad
    w(r)\coloneqq\begin{cases}0,&r\leq\rin,\\ h_{\mathrm{w}}\!\bigl(t(r)\bigr),&\rin<r<\rout,\\ 1,&r\geq\rout,\end{cases}
\end{equation}
where $t(r)=(r-\rin)/(\rout-\rin)$.
Both are $C^3$ on $\R_{>0}$.
The \emph{clamped displacement} is
\begin{equation}\label{eq:rvec-tilde}
    \widetilde{\br}_{ij}\;\coloneqq\;\rtilde(r_{ij})\,\brh_{ij},\qquad \|\widetilde\br_{ij}\|=\rtilde(r_{ij}),
\end{equation}
which preserves direction exactly and replaces $r_{ij}$ by $\rtilde(r_{ij})$ in scalar magnitude.

% =============================================================
\section{Source-Freeze Gate}\label{sec:sfpg}

\paragraph{Per-node gate.}
For each atom $j$, let $\Nb_{\mathrm{out}}(j)\coloneqq\{i:r_{ji}<\rcut,\,i\neq
    j\}$ denote the set of forward neighbours. Define the \emph{source-freeze gate}
\begin{equation}\label{eq:sfpg}
    \eta_j \;\coloneqq\; \prod_{i\in\Nb_{\mathrm{out}}(j)} w\bigl(r_{ji}\bigr) \;\in\;[0,1].
\end{equation}

\begin{lemma}\label{lem:eta}
    $\eta_j$ is a $C^3$ function of $\{\bR_i\}_{i=1}^N$ on the open set $\{\bR_i\neq\bR_j\}$, and
    \begin{equation}\label{eq:eta-zero}
        \exists\,i\in\Nb_{\mathrm{out}}(j)\ \text{with}\ r_{ji}\leq\rin \;\Longrightarrow\; \eta_j=0.
    \end{equation}
\end{lemma}
\begin{proof}
    A finite product of $C^3$ functions is $C^3$.
    The implication follows from $w(r)=0$ for $r\leq\rin$.
\end{proof}

\paragraph{Smooth degree.}
For each atom $i$, define
\begin{equation}\label{eq:deg}
    d_i \;\coloneqq\; \sum_{j\in\Nb_{\mathrm{out}}(i)} s\bigl(r_{ij}\bigr)^2,
    \qquad
    n_i \;\coloneqq\; \bigl(d_i+\varepsilon\bigr)^{-1/2},
\end{equation}
where $\varepsilon>0$ is a fixed regulariser.
Squaring $s$ makes $d_i\in C^6$.

% =============================================================
\section{Truncation to a Reduced Local Layout}\label{sec:trunc}

For computational reasons DPA4 carries, inside the local frame, only those
$(l,m)$ pairs with $|m|\leq\mmax$ for a fixed $\mmax\in\{0,\dots,\lmax\}$.
Define the index set and reduced dimension
\begin{equation}\label{eq:Imm}
    \Idx_{\mmax}\;\coloneqq\;\bigl\{(l,m):0\leq l\leq\lmax,\ |m|\leq\min(l,\mmax)\bigr\},\qquad
    D_m\;\coloneqq\;|\Idx_{\mmax}|.
\end{equation}
Let $\bP_{\mmax}:\R^{D}\to\R^{D_m}$ be the orthogonal projector onto $\Idx_{\mmax}$ and $\bP_{\mmax}^\top$ its adjoint (inclusion).
The truncated edge rotation is
\begin{equation}
    \bD_{ij}^{\le\mmax}\;\coloneqq\;\bP_{\mmax}\bD_{ij}\;\in\;\R^{D_m\times D},
\end{equation}
and the truncated round-trip $\bD_{ij}^{\top}\bP_{\mmax}^\top\bP_{\mmax}\bD_{ij}$ is the orthogonal projector onto the span of $\{\bD_{ij}\,\Idx_{\mmax}\}$ in $V_{\le\lmax}$, hence is not the identity when $\mmax<\lmax$.
To compensate the per-degree norm consumed by truncation, define the diagonal operator $K:\R^{D}\to\R^{D}$ acting on the $(l,m)$ basis element by
\begin{equation}\label{eq:rescale}
    K_{\iota(l,m),\iota(l,m)}\;\coloneqq\;\rho_l,\qquad
    \rho_l\;\coloneqq\;\sqrt{\frac{2l+1}{2\min(l,\mmax)+1}}.
\end{equation}
Two points about $K$:
\begin{itemize}[leftmargin=2em]
    \item It is \emph{not} a single scalar. The value $\rho_l$ depends on the degree $l$
          but is independent of the order $m$, so $K$ acts as the constant $\rho_l$ on
          every row of the degree-$l$ block. Concretely, as a $D\times D$ diagonal
          matrix,
          \begin{equation}\label{eq:K-block}
              K \;=\; \diag\!\bigl(\rho_0\,I_{1},\;\rho_1\,I_{3},\;\rho_2\,I_{5},\dots,\rho_{\lmax}\,I_{2\lmax+1}\bigr),
          \end{equation}
          where $I_{2l+1}$ is the identity on the $(2l+1)$-dimensional block of degree $l$.
          The total number of diagonal entries is $\sum_{l=0}^{\lmax}(2l+1)=D$.
    \item $\rho_l=1$ for every $l\le\mmax$ (no truncation loss), and $\rho_l>1$ for $l>\mmax$ (compensating for the discarded $|m|>\mmax$ coefficients of that degree).
\end{itemize}
The composition $K\,\bD_{ij}^{\top}\bP_{\mmax}^\top$ will be called the \emph{rescaled lift} and is the canonical right inverse used after edge-local processing.

% =============================================================
\section{Equivariant Operators}\label{sec:ops}

We define four classes of parameterised maps that constitute the building
blocks of DPA4. Throughout this section $C,C',F\in\N$ are channel and focus
widths.

\paragraph{Classification by Schur's lemma.}
Before defining the operators we record what they exhaust. For a compact group
$G$ acting on a finite-dimensional real vector space, the real form of Schur's
lemma classifies the space of $G$-equivariant linear maps between isotypic
decompositions:
\[
    \mathrm{Hom}_G(V_a,V_b)=0\ \text{if }V_a\not\cong V_b,\qquad
    \mathrm{End}_G(V_a)\in\{\R,\C,\mathbb{H}\}\ \text{according to the Frobenius--Schur type of }V_a.
\]
Applied to $V_{\le\lmax}\otimes\R^C$ this gives the following two universal
classifications, which are the entire content of the linear theory below.

\emph{Global $\SO(3)$.} Each $V_l$ is a real-type irreducible with $\mathrm{End}_{\SO(3)}(V_l)=\R$, so
\begin{equation}\label{eq:hom-so3}
    \mathrm{Hom}_{\SO(3)}\bigl(V_{\le\lmax}\otimes\R^C,\,V_{\le\lmax}\otimes\R^{C'}\bigr)\;\cong\;\bigoplus_{l=0}^{\lmax}\R^{C\times C'},
\end{equation}
i.e., one independent $C\times C'$ matrix per degree; cross-$l$ mixing is forbidden.
This is exactly the space realised by $L^{\mathrm{deg}}$ below, of which $L^{\mathrm{ch}}$ is the special case in which all degrees share the same matrix.

\emph{Edge-local $\SO(2)_{\bz}$.} The irreducibles are the trivial 1-d rep (each $(l,0)$ line, real type) and the planar-rotation reps $\C_m$ for $m\ge 1$ (complex type, $\mathrm{End}_{\SO(2)_{\bz}}\cong\C$):
    \begin{equation}\label{eq:hom-so2}
        \mathrm{Hom}_{\SO(2)_{\bz}}\bigl(V_{\le\lmax}\otimes\R^C,\,V_{\le\lmax}\otimes\R^{C'}\bigr)\;\cong\;
        \bigoplus_{l,l'}\R^{C\times C'}\;\oplus\;\bigoplus_{m=1}^{\mmax}\!\bigoplus_{\substack{l,l'\\\geq m}}\C^{C\times C'},
    \end{equation}
    where the $m=0$ block allows arbitrary cross-$l$ mixing (trivial rep, $\mathrm{End}=\R$), each fixed-$m>0$ block is governed by complex multiplication realised by the pair $(\bU^{(l,l',m)},\bV^{(l,l',m)})$, and cross-$m$ mixing is forbidden.
    This is exactly the space realised by $L^{\mathrm{SO2}}$ below.

    \paragraph{Beyond linear.}
    The architectural freedom not covered by \eqref{eq:hom-so3}--\eqref{eq:hom-so2}
    consists of: (i) multiplication by an $\SO(3)$-invariant scalar function of
    geometry — used in DPA4 as the per-degree radial profile $\brho_l(r)$, the
    smooth cutoff $s(r)$, and the source-freeze gate $\eta_j$; (ii) gated
    nonlinearities of the form
$\bh_{(l,m)}\mapsto\bh_{(l,m)}\cdot\sigma(\text{invariant})$, used in
$\Gamma_{\psi,\bG}$ below; (iii) bilinear Clebsch--Gordan tensor products
$V_{l_1}\otimes V_{l_2}\to V_{l_3}$ allowed by the SO(3) selection rule
$|l_1-l_2|\le l_3\le l_1+l_2$ (cost $\mathcal{O}(L^6)$ in the truncation
$L=\lmax$), \emph{not} used by DPA4; (iv) their $\SO(2)_{\bz}$ analogue
$\C_{m_1}\otimes\C_{m_2}\to\C_{m_1\pm m_2}$, \emph{not} used by DPA4; (v)
    permutation-equivariant set-aggregation over neighbours, weighted by invariant
    scalars, used in both the envelope-weighted sum and the envelope-gated softmax
    attention. DPA4 adopts (i), (ii), (v) and deliberately omits (iii)--(iv);
    avoiding the cubic-in-$L$ cost of (iii) is the algorithmic motivation for
    working in the local frame.

    \paragraph{Channel mixing.}
    Given $\bW\in\R^{C\times C'}$, the channel-mixing operator
$L^{\mathrm{ch}}_{\bW}:V_{\le\lmax}\otimes\R^C\to V_{\le\lmax}\otimes\R^{C'}$
    is
    \begin{equation}
        \bigl(L^{\mathrm{ch}}_{\bW}\bh\bigr)_{\iota(l,m),c'}\;=\;\sum_{c=1}^{C}\bW_{c,c'}\,\bh_{\iota(l,m),c},
    \end{equation}
    which commutes with $\bD(R)$ for every $R\in\SO(3)$ since it does not mix the $\iota$-axis.

    \paragraph{Degree-wise mixing.}
    Given $\bW=(\bW^{(0)},\dots,\bW^{(\lmax)})\in\prod_{l=0}^{\lmax}\R^{C\times
    C'}$, the degree-wise mixing operator $L^{\mathrm{deg}}_{\bW}$ acts
    independently on each $l$-block:
    \begin{equation}
        \bigl(L^{\mathrm{deg}}_{\bW}\bh\bigr)_{\iota(l,m),c'}\;=\;\sum_{c=1}^{C}\bW^{(l)}_{c,c'}\,\bh_{\iota(l,m),c}.
    \end{equation}
    \begin{proposition}\label{prop:so3lin}
    For every $R\in\SO(3)$, $L^{\mathrm{deg}}_{\bW}\bD(R)=\bD(R)L^{\mathrm{deg}}_{\bW}$.
    The same statement holds, a fortiori, for the channel-mixing special case $L^{\mathrm{ch}}_{\bW}$.
    \end{proposition}
    \begin{proof}
$\bW^{(l)}$ acts only on the channel axis of the degree-$l$ block, so for each $l$ the map commutes with $D^l(R)$, hence with $\bD(R)$ overall.
    \end{proof}

    \paragraph{$\SO(2)$-equivariant mixing.}
    For each pair of degrees $0\le l,l'\le\lmax$, let $\bA^{(l,l',0)}\in\R^{C\times
    C'}$ be an unconstrained linear map (the $m=0$ cross-$l$ mixers). For each
    triple $(l,l',m)$ with $1\le m\le\min(l,l',\mmax)$, let
$\bU^{(l,l',m)},\bV^{(l,l',m)}\in\R^{C\times C'}$ (the $m$-stratum mixers
    between degrees $l$ and $l'$). Optionally, let $\bm{b}_0\in\R^{C'}$ be an
    additive bias acting on the $l=0$ slot only (no bias is allowed on $l\ge 1$,
    since constant shifts of non-trivial irreducibles would break
$\SO(2)_{\bz}$-equivariance). Pack the whole collection as the parameter
$\bW=(\bA,\bU,\bV,\bm{b}_0)$. Define $L^{\mathrm{SO2}}_{\bW}$ on the truncated
    layout by
    \begin{align}
        \bigl(L^{\mathrm{SO2}}_{\bW}\bx\bigr)_{(l,0),c'}
         & = \sum_{l'=0}^{\lmax}\sum_{c=1}^{C}A^{(l,l',0)}_{c,c'}\,\bx_{(l',0),c}\;+\;\bm{b}_{0,c'}\,\delta_{l,0},\label{eq:so2-m0} \\
        \begin{pmatrix}\bigl(L^{\mathrm{SO2}}_{\bW}\bx\bigr)_{(l,-m),c'} \\
            \bigl(L^{\mathrm{SO2}}_{\bW}\bx\bigr)_{(l,+m),c'}\end{pmatrix}
         & = \sum_{l'\geq m}\sum_{c=1}^{C}\begin{pmatrix} U^{(l,l',m)}_{c,c'} & -V^{(l,l',m)}_{c,c'} \\
                                                          V^{(l,l',m)}_{c,c'} & U^{(l,l',m)}_{c,c'}\end{pmatrix}
        \begin{pmatrix}\bx_{(l',-m),c}\\ \bx_{(l',+m),c}\end{pmatrix},\label{eq:so2-mn}
    \end{align}
    where the parameter $\bA^{(l,l',0)}$ allows arbitrary mixing across $l,l'$ within the $m=0$ slice, while in \eqref{eq:so2-mn} mixing within a fixed $|m|$-stratum is restricted to the SO(2)-equivariant $2\times 2$ block-coupling structure.
    Setting $\bm{b}_0\equiv 0$ yields the bias-free variant of $L^{\mathrm{SO2}}$.

    \begin{theorem}\label{thm:so2-eq}
    For every $\theta\in[0,2\pi)$, $L^{\mathrm{SO2}}_{\bW}\,\bD(R_{\theta\bz})=\bD(R_{\theta\bz})\,L^{\mathrm{SO2}}_{\bW}$, where $R_{\theta\bz}\in\SO(2)_{\bz}$ is the rotation by $\theta$ about $\bz$.
    \end{theorem}
    \begin{proof}
    For $m=0$ the action of $\bD(R_{\theta\bz})$ is trivial on every $(l,0)$ subspace, so any linear map among $m=0$ coefficients commutes with it.
    For $m>0$ the $2\times 2$ block in \eqref{eq:so2-mn} is the real realisation of the complex scalar $U+iV$ acting by multiplication on $\C_m$.
    Since $\bD(R_{\theta\bz})$ acts on $\C_m$ by multiplication by $e^{im\theta}$, and complex multiplication is commutative, the operators commute.
    \end{proof}

    \paragraph{Equivariant RMS normalisation.}
    For $\bh\in V_{\le\lmax}\otimes\R^C$, write the scalar mean $\bar h\coloneqq
C^{-1}\sum_{c}\bh_{\iota(0,0),c}$ and the degree-balanced second moment
    \begin{equation}\label{eq:sigma2}
        \sigma^2(\bh) \;\coloneqq\; \sum_{l=0}^{\lmax}\sum_{m=-l}^{+l}\sum_{c=1}^{C}\frac{1}{(2l+1)(\lmax+1)C}\,\bigl(\bh_{\iota(l,m),c}-\delta_{l,0}\,\bar h\bigr)^{2}.
    \end{equation}
    Given per-degree scales $\gamma=(\gamma_0,\dots,\gamma_{\lmax})\in\R^{\lmax+1}$ and an $l=0$ bias $\beta\in\R^{C}$, the \emph{equivariant RMS normalisation} is
    \begin{equation}\label{eq:rmsnorm}
        \bigl(N_{\gamma,\beta}\bh\bigr)_{\iota(l,m),c}\;=\;\gamma_l\cdot\frac{\bh_{\iota(l,m),c}-\delta_{l,0}\,\bar h}{\sqrt{\sigma^2(\bh)+\varepsilon}} \;+\;\delta_{l,0}\,\beta_c.
    \end{equation}

    \begin{proposition}\label{prop:rms-eq}
$N_{\gamma,\beta}$ commutes with $\bD(R)$ for every $R\in\SO(3)$.
    \end{proposition}
    \begin{proof}
    The numerator of \eqref{eq:rmsnorm} is $\SO(3)$-equivariant (subtraction of a scalar from the $l=0$ slice and identity elsewhere).
    The denominator $\sigma^2$ is invariant since each $l$-block transforms by an orthogonal $D^l(R)$.
    The per-degree scaling commutes with $D^l(R)$ for every $l$.
    \end{proof}

    \paragraph{Gated nonlinearity.}
    Let $\psi:\R\to\R$ be a smooth scalar nonlinearity (e.g.\
$\psi(t)=t\,\sigma(t)$, SiLU, where $\sigma(t)=(1+e^{-t})^{-1}$). Given a
    per-degree projection $\bG=(\bG^{(1)},\dots,\bG^{(\lmax)})$ with
$\bG^{(l)}\in\R^{C\times C}$, define
$\Gamma_{\psi,\bG}:V_{\le\lmax}\otimes\R^C\to V_{\le\lmax}\otimes\R^C$ by
    \begin{equation}\label{eq:gated}
        \bigl(\Gamma_{\psi,\bG}\bh\bigr)_{\iota(l,m),c}\;=\;\begin{cases}
            \psi\!\bigl(\bh_{\iota(0,0),c}\bigr),                                                      & l=0,     \\[2pt]
            \bh_{\iota(l,m),c}\cdot\sigma\!\Bigl(\sum_{c'}\bG^{(l)}_{c',c}\,\bh_{\iota(0,0),c'}\Bigr), & l\geq 1.
        \end{cases}
    \end{equation}

    \begin{proposition}
$\Gamma_{\psi,\bG}$ is $\SO(3)$-equivariant.
    \end{proposition}
    \begin{proof}
    The $l=0$ slice is $\SO(3)$-trivial, so $\psi(\cdot)$ on it is invariant.
    For $l\ge 1$, the gate is a scalar function of the invariant $l=0$ slice and therefore commutes with $D^l(R)$ when applied multiplicatively to a degree-$l$ vector.
    \end{proof}

    \paragraph{Scalar-keyed convex combination of equivariant tensors.}\label{sec:dar}
    A general-purpose construction we will reuse is the following: given several
    equivariant tensors of the same shape, form their convex combination with
    weights computed from their $\SO(3)$-invariant scalar content alone.
    Concretely, let $\bx^{(1)},\dots,\bx^{(S)}\in V_{\le\lmax}\otimes\R^C$ be a
    finite collection of node-level equivariant features, all of the same shape,
    and let $\pi:V_{\le\lmax}\otimes\R^C\to\R^{C'}$ be a fixed $\SO(3)$-invariant
    scalar projection (the canonical choice is $\pi(\bx)\coloneqq\bx|_{l=0}$,
    viewed as a vector in $\R^{C}$ after flattening the trivial degree axis; in
    general $\pi$ may include a subsequent RMS normalisation of the $l=0$ slice).
    Choose a learnable query $\bm{q}\in\R^{C'}$; equivalently, in a
    state-conditioned variant, set $\bm{q}\coloneqq W_q\,\pi(\bx^{\mathrm{cur}})$
    for a linear map $W_q\in\R^{C'\times C'}$ and an auxiliary equivariant tensor
$\bx^{\mathrm{cur}}$ of the same shape. Define
    \begin{equation}\label{eq:dar}
        \mathcal{A}\bigl(\bx^{(1)},\dots,\bx^{(S)};\bm{q}\bigr)
        \;\coloneqq\;
        \sum_{s=1}^{S}\alpha_s\,\bx^{(s)}\;\in\;V_{\le\lmax}\otimes\R^C,
        \qquad
        \alpha_s\;\coloneqq\;\frac{\exp\!\langle\bm{q},\pi(\bx^{(s)})\rangle}{\sum_{s'=1}^{S}\exp\!\langle\bm{q},\pi(\bx^{(s')})\rangle}.
    \end{equation}
    At initialisation, $\bm{q}$ (or, in the state-conditioned variant, $W_q$) is set to zero, so $\alpha_s=1/S$ uniformly and $\mathcal{A}$ reduces to the unweighted mean $S^{-1}\sum_s\bx^{(s)}$.

    \begin{proposition}\label{prop:dar-eq}
    For every $R\in\SO(3)$,
$\mathcal{A}(\bD(R)\bx^{(1)},\dots,\bD(R)\bx^{(S)};\bm{q})\;=\;\bD(R)\,\mathcal{A}(\bx^{(1)},\dots,\bx^{(S)};\bm{q})$.
    \end{proposition}
    \begin{proof}
    By assumption $\pi$ takes values in the trivial representation, so the inner products $\langle\bm{q},\pi(\bx^{(s)})\rangle$ are $\SO(3)$-invariant; hence so are the weights $\alpha_s$.
    The output \eqref{eq:dar} is a finite sum of $\SO(3)$-equivariant tensors weighted by $\SO(3)$-invariant scalars, which itself transforms by $\bD(R)$ as a whole.
    \end{proof}

    % =============================================================
    \section{Per-Atom State and Initial Embeddings}\label{sec:init}

    \paragraph{State.}
    A \emph{state} of the message-passing computation is a tuple
$X=(\bh_1,\dots,\bh_N)\in\bigl(V_{\le\lmax}\otimes\R^C\bigr)^{N}$. The group
$S_N$ acts by permuting the index $i$; $\SO(3)$ acts on each $\bh_i$ in the
    first tensor factor.

    \paragraph{Type embedding.}
    Fix a learnable table $T\in\R^{n_t\times C}$. The initial state $X^{(0)}$ has
    \begin{equation}\label{eq:type}
        \bh_i^{(0)}\bigm|_{l=0,m=0,c}\;=\;T_{Z_i,c},\qquad \bh_i^{(0)}\bigm|_{l\geq 1}\;=\;0.
    \end{equation}
    The padding row of $T$, used by masking, is held at zero.

    \paragraph{Geometric initial embedding (GIE).}
    For $l\geq 1$ and $m'\in\{-l,\dots,+l\}$, the $(l,0)$-column of the inverse
    edge rotation $\bD_{ij}^{\top}$ evaluated at row $\iota(l,m')$ is, up to the
    basis normalisation, the real spherical harmonic $Y_l^{m'}(\brh_{ij})$. Denote
    \begin{equation}\label{eq:Y}
        Y_l^{m'}(\brh_{ij})\;\coloneqq\;\bigl[\bD_{ij}^{\top}\bigr]_{\iota(l,m'),\,\iota(l,0)}.
    \end{equation}
    Let $\Te^{(1)},\Te^{(2)}\in\R^{n_t\times C}$ be learnable species embedding tables and define the \emph{per-edge species embedding}
    \begin{equation}\label{eq:Tedge}
        \Te_{\mathrm{edge}}(i,j) \;\coloneqq\; \Te^{(1)}_{Z_i,:}+\Te^{(2)}_{Z_j,:} \;\in\;\R^{C},
    \end{equation}
    a vector that depends on the species pair $(Z_i,Z_j)$ only.
    We treat $\Te_{\mathrm{edge}}(i,j)$ as a vector of length $C$ broadcast uniformly across the degree axis to a tensor in $\R^{(\lmax+1)\times C}$ when added to radial features.

    Let $\Pi^{\mathrm{rad}}:\R^{n_r}\to\R^{(\lmax+1)\times C}$ be a bias-free
    multilayer perceptron with SiLU activations ($\mathrm{SiLU}(t)=t\,\sigma(t)$);
    being bias-free ensures $\Pi^{\mathrm{rad}}(0)=0$, so masked or out-of-cutoff
    edges (for which $\bphi=0$) contribute nothing. The per-edge per-degree radial
    features are
    \begin{equation}\label{eq:rho-edge}
        \brho_{ij} \;\coloneqq\; \Pi^{\mathrm{rad}}\!\bigl(\bphi(\rtilde_{ij})\bigr)+\Te_{\mathrm{edge}}(i,j) \;\in\;\R^{(\lmax+1)\times C},
    \end{equation}
    with $\Te_{\mathrm{edge}}(i,j)\in\R^{C}$ broadcast to all $(\lmax+1)$ degree slots.
    The \emph{geometric initial embedding} adds to the state
    \begin{equation}\label{eq:gie}
        \bh_i^{(0)}\bigm|_{\iota(l,m'),c} \;\mathrel{+}=\; n_i\!\sum_{j:r_{ij}<\rcut}\eta_j\,Y_l^{m'}(\brh_{ij})\,\brho_{ij,l,c},\qquad l\geq 1.
    \end{equation}

    \begin{proposition}\label{prop:gie-eq}
    For every $R\in\SO(3)$, the map $\{\bR_i\}\mapsto\{\bh_i^{(0)}\}$ defined by \eqref{eq:type}-\eqref{eq:gie} satisfies $\bh_i^{(0)}(\{R\bR_j\})=\bD(R)\,\bh_i^{(0)}(\{\bR_j\})$.
    \end{proposition}
    \begin{proof}
    The $l=0$ part is independent of positions.
    The $l\ge 1$ part is a finite sum of products of (a) the scalar $\eta_j\,n_i\,\brho_{ij,l,c}$, which depends only on rotation-invariant data (distances and species), and (b) the degree-$l$ harmonic $Y_l^{m'}(\brh_{ij})$, whose vector $(Y_l^{m'}(\brh))_{m'}$ transforms by $D^l(R)$ under $\brh\mapsto R\brh$.
    \end{proof}

    \paragraph{Environment FiLM (optional).}
    FiLM (\emph{Feature-wise Linear Modulation}, Perez \emph{et al.}, 2018) is an
    affine conditioning $\bx\mapsto\gamma(\bm{c})\odot\bx+\beta(\bm{c})$ in which a
    per-channel scale $\gamma$ and shift $\beta$ are learned functions of a
    conditioning input $\bm{c}$. Here the conditioning is the local geometry seen
    from atom $i$, and the conditioned features are the species-embedded scalar
    slice $\bh_i^{(0)}|_{l=0}$. Define the four-vector $\bu_{ij}\in\R^4$ by
    \begin{equation}\label{eq:utilde}
        \bu_{ij,0}\coloneqq\frac{s(\rtilde_{ij})}{\rtilde_{ij}},\qquad
        \bu_{ij,k}\coloneqq\bu_{ij,0}\,\brh_{ij,k}\ (k=1,2,3).
    \end{equation}
    For a separately parameterised radial--species map $\bg:\R_{\geq 0}\times\{1,\dots,n_t\}^2\to\R^{C_{\mathrm{env}}}$, set
    \begin{equation}
        A_i \;\coloneqq\; n_i\sum_{j:r_{ij}<\rcut}\eta_j\,\bu_{ij}\otimes \bg(\rtilde_{ij};Z_i,Z_j)\;\in\;\R^{4\times C_{\mathrm{env}}}.
    \end{equation}
    The compact descriptor $\mathcal{D}_i\coloneqq A_i^{\top}A_i^{(\cdot,1{:}k)}\in\R^{C_{\mathrm{env}}\times k}$ is $\SO(3)$-invariant because the spatial part of $\bu_{ij}$ transforms by $R\in\SO(3)$ while the temporal component is invariant.
    With learnable log-scalars $\lambda_\alpha,\lambda_\beta\in\R$, weight matrices $W_\alpha,W_\beta$, and an $\SO(3)$-trivial RMS norm $N_0$, the FiLM update on the scalar slice is
    \begin{equation}\label{eq:film}
        \bh_i^{(0)}\bigm|_{l=0}\;\leftarrow\;\bigl[1+\exp(\lambda_\alpha)\tanh\!\bigl(N_0(W_\alpha\,\mathrm{vec}\,\mathcal{D}_i)\bigr)\bigr]\odot\bh_i^{(0)}\bigm|_{l=0}+\exp(\lambda_\beta)\tanh\!\bigl(N_0(W_\beta\,\mathrm{vec}\,\mathcal{D}_i)\bigr).
    \end{equation}
    At initialisation, $\lambda_\alpha=\lambda_\beta=\log(0.01)$, making the FiLM update a near-identity.

    % =============================================================
    \section{Edge Convolution in the Local Frame}\label{sec:so2conv}

    The \emph{edge convolution} is a single $\SO(3)$-equivariant operator
    \begin{equation}\label{eq:Cdef}
        \mathcal{C}_\theta\,:\,\bigl(V_{\le\lmax}\otimes\R^C\bigr)^{N}\;\longrightarrow\;\bigl(V_{\le\lmax}\otimes\R^C\bigr)^{N},
        \qquad
        X\;\longmapsto\;\Delta X,
    \end{equation}
    that performs one round of inter-atomic message passing on the directed neighbour graph: for each atom $i$, the output $\bigl(\mathcal{C}_\theta X\bigr)_i$ aggregates contributions from every $j\in\Nb_{\mathrm{out}}(i)$ in a way that respects the symmetries of \S\ref{sec:setup} and the edge-local frame of \S\ref{sec:frame}.
    The subscript $\theta$ collects every learnable parameter introduced in the remainder of this section.

    The construction proceeds in nine stages. It internally carries an enlarged
    hidden width $H=F\cdot C_f$ with focus count $F\in\N$ and per-focus width
$C_f\in\N$, and contracts back to the descriptor channel width $C$ in the final
    stage. At training time zero, the down-projection of the final stage is
    zero-initialised, so $\mathcal{C}_\theta\equiv 0$ initially.

    \paragraph{Stage 1: pre-mixing.}
    Fix a degree-wise mixing
$L^{\mathrm{deg}}_{\mathrm{pre}}:V_{\le\lmax}\otimes\R^{C}\to
V_{\le\lmax}\otimes\R^{H}$. The hidden representation at every node is
$\bh'_i=L^{\mathrm{deg}}_{\mathrm{pre}}\,\bh_i$.

    \paragraph{Stage 2: local-frame transport.}
    For each directed edge $(i,j)$ with $r_{ij}<\rcut$ define
    \begin{equation}\label{eq:xloc}
        \bx_{ij}\;\coloneqq\;\bD_{ij}^{\le\mmax}\,\bh'_j\;\in\;\R^{D_m\times H}.
    \end{equation}

    \paragraph{Stage 3: per-degree radial modulation.}
    Let $\widetilde\brho_{ij}\coloneqq
L^{\mathrm{ch}}_{\bW_{\mathrm{r}}}\brho_{ij}\in\R^{(\lmax+1)\times H}$ be the
    radial features lifted to the hidden width. The modulation is
    \begin{equation}\label{eq:radmod}
        \bx_{ij,(l,m),c}\;\leftarrow\;\widetilde\rho_{ij,l,c}\cdot\bx_{ij,(l,m),c},\qquad (l,m)\in\Idx_{\mmax},\ c=1,\dots,H.
    \end{equation}

    \paragraph{Stage 4: multi-focus split.}
    View $\R^{H}=\R^{F}\otimes\R^{C_f}$ and write $\bx_{ij}\in\R^{D_m\times F\times
    C_f}$.

    \paragraph{Stage 5: equivariant mixing stack.}
    Let
$\Theta=(L^{\mathrm{SO2}}_1,N_1,\Gamma_1,\dots,L^{\mathrm{SO2}}_S,N_S,\Gamma_S)$
    be a finite sequence of $\SO(2)$ linear maps, equivariant RMS norms, and gated
    activations. By convention, the last layer drops both its pre-norm and its
    activation: $N_S\equiv\mathrm{Id}$ and $\Gamma_S\equiv\mathrm{Id}$. The stack
    action is the composition with residual shortcuts:
    \begin{equation}\label{eq:so2-stack}
        \bx_{ij}\;\leftarrow\;\bx_{ij}+\Lambda_{s}\odot\Gamma_s\!\bigl(L^{\mathrm{SO2}}_s\,N_s(\bx_{ij})\bigr),\qquad s=1,\dots,S,
    \end{equation}
    where $\Lambda_s\in\R^{F\times C_f}$ are learnable per-stream scales (initially $10^{-3}$) and $\odot$ is broadcast multiplication on the focus/channel axes.

    \paragraph{Stage 5a: layer-$0$ bias correction.}
    When the first SO(2) linear map $L^{\mathrm{SO2}}_1$ in the stack carries the
    optional $l=0$ bias $\bm{b}_0$ of \eqref{eq:so2-m0} (here treated as a
    per-focus $\R^{F\times C_f}$-shaped parameter so it can be multiplied
    component-wise by the per-focus, per-channel modulation factor), this bias does
    \emph{not} pass through the radial modulation \eqref{eq:radmod} or the envelope
$s(r_{ij})$. Left as-is, $\bm{b}_0$ would behave as a constant additive
    contribution to the $l=0$ slice of every edge message, including edges whose
    envelope has decayed to zero, breaking the smooth cutoff at $r_{ij}=\rcut$. To
    restore consistency, the implementation \emph{rescales the bias by the same
    modulation factor that multiplies the rest of the layer-$0$ output} and
    subtracts off the original additive bias. Concretely, after applying
$L^{\mathrm{SO2}}_1$ at $s=1$, the $l=0$ slot of $\bx_{ij}$ is updated by
    \begin{equation}\label{eq:bias-corr}
        \bx_{ij,(0,0),f,c}\;\mathrel{+}=\;\bm{b}_{0,f,c}\cdot\bigl(\widetilde\rho_{ij,0,c}\,s(r_{ij})-1\bigr),
    \end{equation}
    which, combined with the bias's nominal addition $+\bm{b}_{0,f,c}$ inside $L^{\mathrm{SO2}}_1$, yields a net contribution of $\bm{b}_{0,f,c}\cdot\widetilde\rho_{ij,0,c}\,s(r_{ij})$ to the layer-$0$ output --- that is, the bias is modulated edge-wise by the same product of radial features and envelope as every other $l=0$ channel, and vanishes $C^3$-smoothly at the cutoff.
    The correction is applied only at $s=1$; subsequent SO(2) layers receive no bias re-treatment because their inputs are already in the modulated regime.

    \paragraph{Stage 5b (optional): intra-stack memory.}
    The plain stack \eqref{eq:so2-stack} can be augmented by the scalar-keyed
    convex combination $\mathcal{A}$ of \eqref{eq:dar}, wired inside the SO(2)
    layer loop so that each layer sees an attention-weighted combination of the
    running history of all earlier intermediate states of the same SO(2) stack.
    Concretely, define the history of equivariant snapshots recursively by
$\mathcal{H}_0=(\bx_{ij}^{[0]})$ with $\bx_{ij}^{[0]}\coloneqq\bx_{ij}$ (the
    input to the stack), and for $s=1,\dots,S$ set
    \begin{equation}\label{eq:so2-depth-attn}
        \bx_{ij}^{(s,\mathrm{in})}\;\coloneqq\;\mathcal{A}\!\bigl(\mathcal{H}_{s-1};\bm{q}^{(s)}\bigr),
    \end{equation}
    where the scalar projection $\pi$ is the canonical one $\bx\mapsto\bx|_{l=0}$, the invariant content of which is well-defined under the residual $\SO(2)_{\bz}$ symmetry of the local frame, and $\bm{q}^{(s)}\in\R^{F\cdot C_f}$ is a per-layer query (either learnable or computed from $\bx_{ij}^{[s-1]}|_{l=0}$ via a per-layer linear map).
    Run the standard layer-$s$ update of \eqref{eq:so2-stack} on $\bx_{ij}^{(s,\mathrm{in})}$, denote its output by $\bx_{ij}^{[s]}$ and the layer increment by $\delta_s\coloneqq\bx_{ij}^{[s]}-\bx_{ij}^{(s,\mathrm{in})}$, and update the history
    \begin{equation}\label{eq:so2-history}
        \mathcal{H}_s\;\coloneqq\;\mathcal{H}_{s-1}\cup\{\delta_s\},
    \end{equation}
    so that layer $s+1$ has access to the initial input and to every per-layer increment produced so far.
$\SO(2)_{\bz}$-equivariance follows from Proposition~\ref{prop:dar-eq} applied to the residual $\SO(2)_{\bz}$ action together with the layer-$s$ equivariance of \eqref{eq:so2-stack}.
    At initialisation $\bm{q}^{(s)}=0$, so the attention weights are uniform across $\mathcal{H}_{s-1}$ and \eqref{eq:so2-depth-attn} reduces to the unweighted mean of the history; the model learns layer-specific attention patterns from the loss.
    This component is optional; when disabled, the stack reverts to the plain composition \eqref{eq:so2-stack}.

    \paragraph{Stage 6: cross-focus competition.}
    Let $\bx^{(0)}_{ij}\in\R^{F\times C_f}$ denote the $(l,m)=(0,0)$ slice of
$\bx_{ij}$ at the entry of the stack \eqref{eq:so2-stack}. For a learnable
    parameter $W^{\mathrm{cf}}\in\R^{C_f\times F}$, temperature $\tau>0$, and
    label-smoothing parameter $\epsilon\in[0,1)$, define
    \begin{equation}\label{eq:focuscomp}
        \zeta_{ij,f}=\frac{\exp\!\bigl(\tau^{-1}\sum_{c}W^{\mathrm{cf}}_{c,f}\,N_0(\bx^{(0)}_{ij})_{f,c}\bigr)}{\sum_{f'}\exp\!\bigl(\tau^{-1}\sum_{c}W^{\mathrm{cf}}_{c,f'}\,N_0(\bx^{(0)}_{ij})_{f',c}\bigr)},
        \quad
        \alpha_{ij,f}=(1-\epsilon)\zeta_{ij,f}+\epsilon/F,
    \end{equation}
    and reweight $\bx_{ij}\leftarrow\alpha_{ij}\odot\bx_{ij}$ (with $\alpha$ broadcast across the $(l,m)$ and $c$ axes).
    Because $\bx^{(0)}_{ij}$ is $\SO(2)_{\bz}$-invariant, this multiplication does not break the equivariance of Theorem~\ref{thm:so2-eq}.

    \paragraph{Stage 7: global lift and rescale.}
    The edge message is
    \begin{equation}\label{eq:edge-msg}
        \bm{m}_{ij}\;\coloneqq\;K\,\bigl(\bD_{ij}^{\le\mmax}\bigr)^{\top}\,\bx_{ij}\;\in\;\R^{D\times H},
    \end{equation}
    with $K$ as in \eqref{eq:rescale}.

    \paragraph{Stage 8: aggregation over neighbours (choice of A or B).}
    The edge messages $\{\bm{m}_{ij}\}_{j\in\Nb_{\mathrm{out}}(i)}$ are combined
    into a single per-atom update $(\Delta\bh')_i$ by exactly one of the two
    mutually exclusive variants below. The two variants share the same input, the
    same output type $(\Delta\bh')_i\in\R^{D\times H}$, and feed into the same
    downstream Stage 9; they differ only in how they weight the contributions of
    individual neighbours. Variant A is the default; Variant B is enabled when an
    attention-head count $H_a\in\N$ is configured.

    \subparagraph{Variant A --- envelope-weighted sum.}
    The plain aggregation produces
    \begin{equation}\label{eq:agg-plain}
        \bigl(\Delta\bh'\bigr)_i\;\coloneqq\;n_i\!\sum_{j:r_{ij}<\rcut} s(r_{ij})\,\eta_j\,\bm{m}_{ij}\;\in\;\R^{D\times H}.
    \end{equation}
    The combined edge weight $s(r_{ij})\,\eta_j$ is an $\SO(3)$-invariant scalar that acts uniformly on all $(l,m)$ components of $\bm{m}_{ij}$; it is therefore an additional ``scalar $\times$ equivariant'' multiplication on top of the per-degree radial modulation \eqref{eq:radmod}, and introduces no new equivariant edge feature.
    In total, the only edge-side scalars that ever multiply the equivariant message between Stage 2 and aggregation are the per-degree radial features $\brho_{ij,l}$ (in-frame, per-degree), the truncation rescale $\rho_l$ (a fixed factor, no parameters), and the envelope/gate product $s(r_{ij})\,\eta_j$ (global per-edge); all four are $\SO(3)$-invariant.

    \subparagraph{Variant B --- envelope-gated multi-head attention.}
    The attention path replaces the plain envelope-weighted sum
    \eqref{eq:agg-plain} by a destination-wise normalised sum whose weights are
    learned scalar functions of $l=0$ invariants. It is built in seven explicit
    substeps.

    Choose an attention-head count $H_a\in\N$ that divides the per-focus width
$C_f$, and let $d_a\coloneqq C_f/H_a$. The total hidden width thus factorises
    as $H=F\cdot C_f=F\cdot H_a\cdot d_a$. Let
    \[
    Q,\,K_{\mathrm{att}}\in\R^{F\times C_f\times C_f},\quad
    W^{\mathrm{rb}}\in\R^{C_f\times F\times H_a},\quad
    W^{\mathrm{og}}\in\R^{C_f\times F\times H_a},\quad
    \zeta_a\in\R^{F\times H_a}
\]
be learnable parameters; $Q,K_{\mathrm{att}}$ act as focus-wise linear maps
(one independent $C_f\times C_f$ matrix per focus).

\emph{(B1) Queries and keys.}
The $\SO(3)$-invariant $l=0$ slice $\bh_n|_{l=0}\in\R^{H}$ of every node feature is reshaped on the channel axis as $\R^{F\times C_f}$ and normalised by a focus-wise scalar RMS norm $N_0$.
Define
\begin{equation}\label{eq:qk}
    \bq_n \;\coloneqq\; Q\,N_0\bigl(\bh_n|_{l=0}\bigr),\qquad
    \bk_n \;\coloneqq\; K_{\mathrm{att}}\,N_0\bigl(\bh_n|_{l=0}\bigr),\qquad n=1,\dots,N,
\end{equation}
both in $\R^{F\times C_f}$, then split the $C_f$-axis into $H_a$ heads of width $d_a$:
\begin{equation}\label{eq:qk-split}
    \bq_n\;\in\;\R^{F\times H_a\times d_a},\qquad \bk_n\;\in\;\R^{F\times H_a\times d_a}.
\end{equation}

\emph{(B2) Edge logits with radial bias.}
For each edge $(i,j)$, focus $f$, and head $h$, the scaled dot product is augmented with a per-edge radial bias derived from the $l=0$ slice of the lifted radial features $\widetilde\rho_{ij,0,c}$:
\begin{equation}\label{eq:logit}
    \ell_{ij}^{(f,h)} \;\coloneqq\; \frac{\langle \bq_{i,f,h,:},\,\bk_{j,f,h,:}\rangle}{\sqrt{d_a}}
    \;+\; \sum_{c=1}^{C_f}W^{\mathrm{rb}}_{c,f,h}\,\widetilde\rho_{ij,0,c}.
\end{equation}

\emph{(B3) Destination-wise stable softmax with envelope-and-gate weighting.}
For every destination $i$ and every $(f,h)$, set $\mu_i^{(f,h)}\coloneqq\max_{k:r_{ki}<\rcut}\ell_{ki}^{(f,h)}$, with the convention that edges having $s(r_{ki})^2\,\eta_k=0$ contribute $-\infty$ before the maximum so they are excluded from the normalisation.
The attention weight is
\begin{equation}\label{eq:attw}
    \alpha_{ij}^{(f,h)} \;\coloneqq\;
    \frac{s(r_{ij})^2\,\eta_j\,\exp\!\bigl(\ell_{ij}^{(f,h)}-\mu_i^{(f,h)}\bigr)}
    {\softplus(\zeta_a^{(f,h)})\,\exp\!\bigl(-\mu_i^{(f,h)}\bigr)\;+\;\displaystyle\sum_{k:r_{ki}<\rcut}s(r_{ki})^2\,\eta_k\,\exp\!\bigl(\ell_{ki}^{(f,h)}-\mu_i^{(f,h)}\bigr)}.
\end{equation}
Two mechanisms work together to make the weights well-defined and smooth at the cutoff:
\begin{enumerate}[leftmargin=2em,label=(\roman*)]
    \item the numerator factor $s(r_{ij})^2$ drives $\alpha_{ij}^{(f,h)}\to 0$ as
          $r_{ij}\to\rcut^{-}$, $C^3$-smoothly together with its first three derivatives,
          by the regularity of the cutoff envelope (Lemma~\ref{lem:envelope});
    \item the term $\softplus(\zeta_a^{(f,h)})\exp(-\mu_i^{(f,h)})$ in the denominator is
          a strictly positive regulariser that keeps the denominator bounded below by
          $\softplus(\zeta_a^{(f,h)})\exp(-\mu_i^{(f,h)})>0$ even when every incident
          edge of atom $i$ is silenced (all $s(r_{ki})\to 0$ or $\eta_k\to 0$), so the
          limit $\alpha_{ij}^{(f,h)}\to 0$ is well-defined and free of $0/0$
          indeterminacies.
\end{enumerate}

\emph{(B4) Head split of the message.}
The full-frame message $\bm{m}_{ij}\in\R^{D\times H}$ of \eqref{eq:edge-msg} is reshaped on its channel axis as
\begin{equation}\label{eq:m-head-split}
    \bm{m}_{ij}\;\in\;\R^{D\times F\times H_a\times d_a},
\end{equation}
yielding per-(focus, head) value slices $\bm{m}_{ij}^{(f,h)}\in\R^{D\times d_a}$.

\emph{(B5) Per-(focus, head) attention aggregation.}
\begin{equation}\label{eq:attagg}
    \bA_i^{(f,h)} \;\coloneqq\; \sum_{j:r_{ij}<\rcut}\alpha_{ij}^{(f,h)}\,\bm{m}_{ij}^{(f,h)} \;\in\;\R^{D\times d_a}.
\end{equation}
No $1/\sqrt{\deg_i}$ rescale appears: the softmax denominator in \eqref{eq:attw} already provides the per-destination normalisation that $n_i$ would have supplied in the plain path \eqref{eq:agg-plain}.

\emph{(B6) Destination-side output gate.}
Computed from atom $i$'s $\SO(3)$-invariant $l=0$ slice via the same focus-wise scalar RMSNorm $N_0$,
\begin{equation}\label{eq:outgate}
    G_i^{(f,h)}\;\coloneqq\;\sigma\!\Bigl(\sum_{c=1}^{C_f}W^{\mathrm{og}}_{c,f,h}\,N_0\!\bigl(\bh_i|_{l=0}\bigr)_{f,c}\Bigr)\;\in\;(0,1).
\end{equation}
This gate multiplies each (focus, head) slot of the aggregated value:
\begin{equation}\label{eq:gated-agg}
    \bigl(\Delta\bh'\bigr)_i^{(f,h)} \;\coloneqq\; G_i^{(f,h)}\,\bA_i^{(f,h)} \;\in\;\R^{D\times d_a}.
\end{equation}

\emph{(B7) Merge of (focus, head, head-dim) into the hidden width.}
The four axes $(D,F,H_a,d_a)$ are concatenated back to $(D,H)$ with $H=F\cdot H_a\cdot d_a$:
\begin{equation}\label{eq:agg-attn}
    \bigl(\Delta\bh'\bigr)_i \;\in\;\R^{D\times H}.
\end{equation}
This $(\Delta\bh')_i$ is the analogue of \eqref{eq:agg-plain} on the attention path; the subsequent post-mixing step $L^{\mathrm{deg}}_{\mathrm{post}}:\R^{D\times H}\to\R^{D\times C}$ then projects the result back to the descriptor channel width $C$.

In summary, the construction is
\[
    \boxed{\;
        \bigl(\Delta\bh'\bigr)_i^{(f,h)}\;=\;G_i^{(f,h)}\;\sum_{j:r_{ij}<\rcut}\;\alpha_{ij}^{(f,h)}\,\bm{m}_{ij}^{(f,h)}\;
    }\quad\text{(with $\alpha_{ij}^{(f,h)}$ from \eqref{eq:attw} and $G_i^{(f,h)}$ from \eqref{eq:outgate})},
\]
followed by the focus-and-head merge \eqref{eq:agg-attn}. Every scalar
appearing here ($\alpha$, $G$, $s$, $\eta$, $\zeta_a$, $\widetilde\rho_{ij,0}$)
is $\SO(3)$-invariant; the only equivariant object on the right-hand side is
$\bm{m}_{ij}^{(f,h)}$.

\paragraph{Stage 9: post-mixing.}
A final degree-wise map
$L^{\mathrm{deg}}_{\mathrm{post}}:V_{\le\lmax}\otimes\R^{H}\to
    V_{\le\lmax}\otimes\R^{C}$, with weights initialised at zero, returns the
update to the descriptor channel width:
\begin{equation}\label{eq:Cpostmix}
    \bigl(\mathcal{C}_\theta X\bigr)_i \;\coloneqq\; L^{\mathrm{deg}}_{\mathrm{post}}\bigl((\Delta\bh')_i\bigr) \;\in\;\R^{D\times C}.
\end{equation}
This completes the definition of the operator $\mathcal{C}_\theta$ promised in \eqref{eq:Cdef}: $\theta$ is the tuple of all parameters introduced in Stages 1--9 (the pre/post degree-wise maps, the per-degree radial map $\widetilde\brho$, the SO(2) stack $\Theta$ with the optional intra-stack memory, the cross-focus weights, the attention parameters when active, and the output gate).

\begin{theorem}\label{thm:conv-eq}
    $\mathcal{C}_{\theta}$ is $\SO(3)$-equivariant: under $\bR_i\mapsto R\bR_i$, the increment $\Delta\bh_i$ transforms as $\bD(R)\,\Delta\bh_i$.
\end{theorem}
\begin{proof}
    Under a global rotation, $\bq_{ij}\mapsto\bq_{ij}\cdot R^{-1}$ (in the sense of the double cover), so $\bD_{ij}\mapsto\bD_{ij}\bD(R)^{-1}$.
    Therefore $\bx_{ij}$ is invariant by \eqref{eq:xloc}.
    The stack \eqref{eq:so2-stack} acts within the local frame and preserves the layout (Theorem~\ref{thm:so2-eq}, Proposition~\ref{prop:rms-eq}, gated activation equivariance).
    The cross-focus weights \eqref{eq:focuscomp} are $\SO(2)_{\bz}$-invariant.
    The lift $K\,(\bD_{ij}^{\le\mmax})^{\top}$ then maps the invariant local features to the global frame as $\bD(R)\,K\,(\bD_{ij}^{\le\mmax})^{\top}\bx_{ij}$ on transformed configurations.
    Aggregation \eqref{eq:agg-plain} or attention is a sum of these messages weighted by $\SO(3)$-invariant scalars ($s$, $\eta$, attention weights), hence inherits the same transformation rule.
    $L^{\mathrm{deg}}_{\mathrm{post}}$ commutes with $\bD(R)$.
\end{proof}

% =============================================================
\section{The Equivariant Feed-Forward Network}\label{sec:effn}

The equivariant feed-forward block
$\mathcal{F}_\theta:V_{\le\lmax}\otimes\R^C\to V_{\le\lmax}\otimes\R^C$ is
parameterised by a tuple $\theta$ collecting two degree-wise linear maps and
the parameters of an equivariant nonlinearity. Three variants of
$\mathcal{F}_\theta$ are admitted, all $\SO(3)$-equivariant by
Propositions~\ref{prop:so3lin} and the equivariance of their respective
nonlinearities. The second linear map of each variant is zero-initialised, so
$\mathcal{F}_\theta\equiv 0$ at training time zero and the residual stream
around $\mathcal{F}_\theta$ starts as the identity.

\paragraph{Variant 1: gated FFN (standard).}
\begin{equation}\label{eq:effn-std}
    \mathcal{F}_\theta^{\mathrm{std}}\,\bh\;\coloneqq\;L^{\mathrm{deg}}_{\theta_2}\,\Gamma_{\psi,\bG}\,L^{\mathrm{deg}}_{\theta_1}\,\bh,
\end{equation}
with $L^{\mathrm{deg}}_{\theta_1}:V_{\le\lmax}\otimes\R^C\to V_{\le\lmax}\otimes\R^{C_h}$ (an up-projection by a factor $C_h\geq C$), $L^{\mathrm{deg}}_{\theta_2}:V_{\le\lmax}\otimes\R^{C_h}\to V_{\le\lmax}\otimes\R^C$ (zero-initialised down-projection), and $\Gamma_{\psi,\bG}$ the gated activation of \eqref{eq:gated}.

When the optional GLU split is enabled, $L^{\mathrm{deg}}_{\theta_1}$ instead
outputs a doubled channel width $2C_h$, and the channel axis of its output is
partitioned as $\R^{2C_h}=\R^{C_h}\oplus\R^{C_h}$ into a \emph{value} block
$\bm{v}$ and a \emph{gate} block $\bm{g}$; the activation is replaced by the
pointwise pair-wise product $\bm{v}\odot\psi(\bm{g})$ on each $(l,m)$ slot,
followed by the $\Gamma_{\psi,\bG}$-style gating on $l\geq 1$ channels from
$l=0$. This is the standard SwiGLU variant when $\psi=\SiLU$.

\paragraph{Variant 2: spherical-grid FFN.}
For applications where channel-wise nonlinearity on equivariant features is not
flexible enough, $\mathcal{F}_\theta$ may instead map the higher-degree slices
to a finite sampling of the sphere $\Sph$, apply a pointwise nonlinearity
there, and project back. Fix a finite set of $A\in\N$ quadrature points
$\{(\theta_a,\phi_a)\}_{a=1}^A\subset\Sph$ with associated weights
$\{w_a\}_{a=1}^A\subset\R_{>0}$. Real spherical harmonics evaluated at the
quadrature points give a pair of linear maps (independent of the configuration)
\begin{equation}\label{eq:s2-grid}
    \mathsf{P}:\;V_{\le\lmax}\to\R^{A},\qquad \mathsf{P}^{\dagger}:\;\R^{A}\to V_{\le\lmax},
\end{equation}
where $(\mathsf{P}\bx)_a=\sum_{l,m}Y_l^m(\theta_a,\phi_a)\,\bx_{(l,m)}$ samples a degree-$l$ function on the chosen grid, and $\mathsf{P}^{\dagger}$ is the corresponding quadrature-weighted adjoint, $((\mathsf{P}^{\dagger}\bm{f}))_{(l,m)}=\sum_{a}w_a\,Y_l^m(\theta_a,\phi_a)\,\bm{f}_a$.
Both maps act channel-wise; both are $\SO(3)$-covariant in the sense that $\mathsf{P}^{\dagger}\mathsf{P}=\Pi_{\le\lmax}$, the orthogonal projector onto $V_{\le\lmax}$, when the quadrature is exact on degree-$2\lmax$ trigonometric polynomials.

The spherical-grid FFN is then
\begin{equation}\label{eq:effn-grid}
    \mathcal{F}_\theta^{\mathrm{grid}}\,\bh
    \;\coloneqq\;
    L^{\mathrm{deg}}_{\theta_2}\;\bigl(\,\bh_0' + \mathsf{P}^{\dagger}\circ \Psi \circ \mathsf{P}\,(\bh - \bh|_{l=0})\,\bigr),
\end{equation}
where:
\begin{itemize}[leftmargin=2em]
    \item $\bh_0'\coloneqq L^{\mathrm{ch}}_W\bigl(\psi\bigl(L^{\mathrm{ch}}_{W'}(\bh|_{l=0})\bigr)\bigr)$ is a scalar SwiGLU/GLU branch applied to the $l=0$ slice (a channel-wise MLP, $\SO(3)$-invariant because $l=0$ is the trivial representation), embedded back into the $l=0$ slot of $V_{\le\lmax}$;
    \item $\Psi:\R^{A\times C_h}\to\R^{A\times C_h}$ is a pointwise channel-mixing nonlinearity at each quadrature point — a small per-point MLP with SiLU activations and bias-free linear layers, identical at every point of the grid;
    \item $L^{\mathrm{deg}}_{\theta_1}$ up-projects to width $C_h$ before $\mathsf{P}$; $L^{\mathrm{deg}}_{\theta_2}$ is the zero-initialised down-projection at the end.
\end{itemize}
Pointwise operations on the grid are $\SO(3)$-equivariant only as a composition with $\mathsf{P}$ and $\mathsf{P}^{\dagger}$: the underlying $\SO(3)$ acts on the grid by permutation/interpolation under the standard quadrature, and the pair $(\mathsf{P}^{\dagger}\circ\,\cdot\,\circ\mathsf{P})$ intertwines this action with the action on $V_{\le\lmax}$.
Hence $\mathcal{F}_\theta^{\mathrm{grid}}$ is $\SO(3)$-equivariant.

\paragraph{Variant 3: spherical-grid SwiGLU activation.}
A hybrid of Variants 1 and 2: the first linear map
$L^{\mathrm{deg}}_{\theta_1}$ produces a doubled hidden width $2C_h$, and a
\emph{spherical-grid SwiGLU} activation $\Gamma^{\mathrm{S2}}$ replaces
$\Gamma_{\psi,\bG}$:
\begin{equation}\label{eq:effn-s2glu}
    \mathcal{F}_\theta^{\mathrm{S2}}\,\bh \;\coloneqq\; L^{\mathrm{deg}}_{\theta_2}\,\Gamma^{\mathrm{S2}}\!\bigl(L^{\mathrm{deg}}_{\theta_1}\,\bh\bigr).
\end{equation}
The action of $\Gamma^{\mathrm{S2}}$ on a hidden tensor $\bm{y}\in V_{\le\lmax}\otimes\R^{2C_h}$ is:
\begin{enumerate}[leftmargin=2em,label=(\alph*)]
    \item split the channel axis $\R^{2C_h}=\R^{C_h}\oplus\R^{C_h}$ into a value block
          $\bm{v}$ and a gate block $\bm{g}$;
    \item on the $l=0$ slice, return the SwiGLU $\bm{v}|_{l=0}\odot\psi(\bm{g}|_{l=0})$
          on the $C_h$-dimensional channel axis (a scalar pointwise operation);
    \item on the $l\geq 1$ slices, send the value block through $\mathsf{P}$ to the
          spherical grid, multiply pointwise on the grid by $\psi$ applied to the grid
          image of the gate block, send back through $\mathsf{P}^{\dagger}$;
    \item further multiply the resulting $l\geq 1$ slices by a sigmoidal scalar gate
          computed from the $l=0$ slice (as in $\Gamma_{\psi,\bG}$ of \eqref{eq:gated}).
\end{enumerate}
This produces a single output of width $C_h$, ready for the down-projection by $L^{\mathrm{deg}}_{\theta_2}$.

\paragraph{Equivariance summary.}
All three variants share the same boundary conditions:
$L^{\mathrm{deg}}_{\theta_2}$ is zero-initialised at training time zero (so
$\mathcal{F}_\theta\equiv 0$ initially), and each variant is
$\SO(3)$-equivariant. For Variant 1 this follows from
Propositions~\ref{prop:so3lin} and the equivariance of $\Gamma_{\psi,\bG}$. For
Variants 2 and 3, the spherical-grid maps $\mathsf{P},\mathsf{P}^{\dagger}$
intertwine the $\SO(3)$ action on $V_{\le\lmax}$ with the standard $\SO(3)$
action on $\R^A$ (rotation of quadrature points), so pointwise operations on
the grid become $\SO(3)$-equivariant when sandwiched between $\mathsf{P}$ and
$\mathsf{P}^{\dagger}$.

% =============================================================
\section{The Interaction Block}\label{sec:block}

The interaction block $\mathcal{B}_\theta:X\mapsto X'$ is the composition of
one edge convolution and $n_{\mathrm{f}}\in\N$ equivariant feed-forward blocks,
each wrapped in a residual shortcut and bracketed by optional equivariant RMS
norms on either side of the residual branch (the "sandwich-norm" pattern). Let
\[
    \bigl(\nu^{\mathrm{c,pre}},\nu^{\mathrm{c,post}},\nu^{\mathrm{f,pre}},\nu^{\mathrm{f,post}}\bigr)\in\{0,1\}^{4}
\]
be a 4-bit configuration choosing, independently for each branch, whether to
apply pre- and post-norms. For each enabled position, instantiate an
equivariant RMS norm $N$ (in the sense of \eqref{eq:rmsnorm}, possibly with its
own learnable parameters); for each disabled position, set the corresponding
$N\equiv\mathrm{Id}$. With these conventions, $\mathcal{B}_\theta$ acts on a
state $X=(\bh_1,\dots,\bh_N)$ by
\begin{equation}\label{eq:block}
    \begin{aligned}
        \bh_i^{(1)}   & = \bh_i \;+\; N^{\mathrm{c,post}}\!\Bigl(\mathcal{C}_\theta\bigl(N^{\mathrm{c,pre}}\bh\bigr)\Bigr)_i,                                                                     \\
        \bh_i^{(u+1)} & = \bh_i^{(u)} \;+\; \lambda_u\odot N^{\mathrm{f,post}}_u\!\Bigl(\mathcal{F}_{\theta_u}\bigl(N^{\mathrm{f,pre}}_u\,\bh_i^{(u)}\bigr)\Bigr),\quad u=1,\dots,n_{\mathrm{f}}, \\
        \bh_i'        & = \bh_i^{(n_{\mathrm{f}}+1)},
    \end{aligned}
\end{equation}
where $\lambda_u\in\R^{C}$ are learnable per-channel scales initialised at $10^{-3}$ and acting by broadcast multiplication on the channel axis of each $(l,m)$ slot.
By Proposition~\ref{prop:rms-eq} and the equivariance of $\mathcal{C}_\theta$ (Theorem~\ref{thm:conv-eq}) and $\mathcal{F}_{\theta_u}$ (\S\ref{sec:effn}), every term of \eqref{eq:block} is $\SO(3)$-equivariant.

The four norm switches collapse to four practical patterns:
\begin{itemize}[leftmargin=2em]
    \item \emph{pre-norm only} ($\nu^{\mathrm{c,pre}}=\nu^{\mathrm{f,pre}}=1$, the rest $0$): the default of \eqref{eq:block}; equivalent to the standard Pre-LN Transformer-style stabilisation, with the residual stream itself unnormalised across blocks.
    \item \emph{post-norm only} ($\nu^{\mathrm{c,post}}=\nu^{\mathrm{f,post}}=1$): each residual contribution is normalised after the (non-linear) update, with the residual stream itself untouched.
    \item \emph{sandwich norm} ($\nu^{\mathrm{c,pre}}=\nu^{\mathrm{c,post}}=\nu^{\mathrm{f,pre}}=\nu^{\mathrm{f,post}}=1$): both pre- and post-norms active, mirroring the bracketed structure used in some modern equivariant Transformers.
    \item \emph{no-norm} (all four $=0$): the bare residual composition $\bh_i'=\bh_i+\mathcal{C}_\theta(\bh)_i+\sum_u\lambda_u\odot\mathcal{F}_{\theta_u}(\bh_i^{(u)})$.
\end{itemize}
At initialisation, the second linear maps of both $\mathcal{C}_\theta$ ($L^{\mathrm{deg}}_{\mathrm{post}}$) and each $\mathcal{F}_{\theta_u}$ (the down-projection of the FFN) are zero (cf.\ \S\ref{sec:so2conv}, \S\ref{sec:effn}), so every additive term in \eqref{eq:block} vanishes and $\mathcal{B}_\theta$ starts as the identity map on the state.

% =============================================================
\section{The DPA4 Energy}\label{sec:energy}

\paragraph{Block stack.}
Fix a number of blocks $B\in\N$ and two non-increasing sequences
$\lmax^{(1)}\geq\dots\geq\lmax^{(B)}$ and $\mmax^{(b)}\leq\lmax^{(b)}$. Block
$b$ operates on the truncation of the state to the subspace
$V_{\le\lmax^{(b)}}\otimes\R^C$ and uses $(\lmax^{(b)},\mmax^{(b)})$ in its
edge convolution. Let $\mathcal{B}_{\theta_b}^{(b)}$ denote block $b$ and
define
\begin{equation}\label{eq:final-state}
    X^{(B)} \;\coloneqq\;\mathcal{B}_{\theta_B}^{(B)}\circ\cdots\circ\mathcal{B}_{\theta_1}^{(1)}\,X^{(0)}.
\end{equation}

\paragraph{Scalar readout.}
Let $\mathcal{O}:\R^{C}\to\R$ be a scalar MLP and define the per-atom learned
energy
\begin{equation}\label{eq:Ei}
    E_i^{\mathrm{NN}}\;\coloneqq\;\mathcal{O}\!\bigl(\bh_i^{(B)}\bigm|_{l=0}\bigr),\qquad
    E^{\mathrm{NN}}\;\coloneqq\;\sum_{i=1}^{N}E_i^{\mathrm{NN}}.
\end{equation}

\paragraph{Optional case FiLM (multi-task only).}
In the multi-task setting, where a single network is trained jointly on $K\geq
    2$ datasets, an optional second FiLM conditioner is inserted into the readout
$\mathcal{O}$. Given a one-hot case index $\bm{c}_t\in\R^{K}$ for the dataset
of frame $t$, learnable matrices
$W_\alpha^{\mathrm{case}},W_\beta^{\mathrm{case}}\in\R^{C\times K}$ produce a
per-case scale and shift $\gamma_t=1+W_\alpha^{\mathrm{case}}\bm{c}_t$,
$\beta_t=W_\beta^{\mathrm{case}}\bm{c}_t$, and the input to the scalar MLP is
modulated as $\bh_i^{(B)}|_{l=0}\;\mapsto\;\gamma_t\odot
    \bh_i^{(B)}|_{l=0}+\beta_t$. This case FiLM modifies only the readout, not the
descriptor; the descriptor $\bh_i^{(B)}$ remains case-independent, and the
rotation-invariance, smoothness, and freezing properties of \S\ref{sec:thms}
are unaffected. In the single-task setting considered throughout this document,
$K=1$ and the case FiLM reduces to the identity.

\paragraph{Analytical ZBL branch.}\label{sec:zbl}
Let $k_e>0$ be the Coulomb constant in the chosen units, $a_0>0$ the Bohr
radius, and define for ordered pairs $(i,j)$ with $r_{ij}>0$
\begin{equation}\label{eq:zbl}
    E_{ij}^{\mathrm{ZBL}}(r)\;\coloneqq\;\frac{k_e\,Z_iZ_j}{r}\,\Phi\!\Bigl(\frac{r}{a_{ij}}\Bigr),\qquad
    a_{ij}\;\coloneqq\;\frac{0.88534\,a_0}{Z_i^{0.23}+Z_j^{0.23}},
\end{equation}
\begin{equation}\label{eq:zbl-Phi}
    \Phi(x)\;\coloneqq\;0.18175\,e^{-3.1998x}+0.50986\,e^{-0.94229x}+0.28022\,e^{-0.4029x}+0.028171\,e^{-0.20162x}.
\end{equation}
The analytical branch is $E^{\mathrm{ZBL}}\coloneqq\tfrac{1}{2}\sum_{i\neq j}E_{ij}^{\mathrm{ZBL}}(r_{ij})$.

\paragraph{Bridged total energy.}
The geometric inputs consumed by the learned branch are the clamped pairs
$\{(\widetilde\br_{ij},\rtilde_{ij})\}$ defined in \eqref{eq:rvec-tilde},
whereas the ZBL branch consumes the true distances $\{r_{ij}\}$. The total
energy is
\begin{equation}\label{eq:Etot}
    E\bigl(Z,R\bigr)\;=\;E^{\mathrm{NN}}\bigl(Z,R;\{\widetilde\br_{ij}\}\bigr)+E^{\mathrm{ZBL}}\bigl(Z,R;\{r_{ij}\}\bigr).
\end{equation}
Forces and virial are the gradients of $E$ with respect to positions and the cell.

% =============================================================
\section{Properties}\label{sec:thms}

\begin{theorem}[Symmetry]\label{thm:symm}
    The map $(Z,R)\mapsto E(Z,R)$ defined by \eqref{eq:Etot} is invariant under translations of $R$, under species-preserving permutations of atom indices, and under $\SO(3)$ acting diagonally on positions.
\end{theorem}
\begin{proof}
    Translational invariance: every operation depends on $R$ only through the differences $\bR_j-\bR_i$.
    Permutation invariance: aggregation over $\Nb_{\mathrm{out}}(i)$ is by unordered sum; all parameters depend on $Z_i$, not on $i$.
    $\SO(3)$-invariance: Proposition~\ref{prop:gie-eq} ensures $X^{(0)}$ transforms as $\bD(R)X^{(0)}$; Theorem~\ref{thm:conv-eq} and Proposition~\ref{prop:rms-eq} (and equivariance of $\mathcal{F}_\theta$ and $\Gamma_{\psi,\bG}$) preserve this through every block, so $X^{(B)}$ transforms by $\bD(R)$.
    $E_i^{\mathrm{NN}}$ depends only on the $l=0$ slice of $\bh_i^{(B)}$, which lies in the trivial representation.
    $E^{\mathrm{ZBL}}$ depends only on scalar distances.
\end{proof}

\begin{theorem}[Smoothness]\label{thm:smooth}
    $E\in C^3$ on the open set $\Omega\coloneqq\{R:\bR_i\neq\bR_j\ \forall i\neq j\}$.
\end{theorem}
\begin{proof}
    $s_p\in C^3$ (Lemma~\ref{lem:envelope}), $\rtilde,w\in C^3$ on $\R_{>0}$ by construction, $\bphi\in C^\infty$ on $\R_{>0}$, the two-chart quaternion is $C^\infty$ on $\Sph$ by \eqref{eq:nlerp}, and hence $\bD_{ij}\in C^\infty$ on $\Omega$.
    Every parametric map introduced in \S\ref{sec:ops}, \S\ref{sec:init}, \S\ref{sec:so2conv}, and \S\ref{sec:effn} is smooth on its domain, and the denominator of the attention weight \eqref{eq:attw} is bounded below by $\softplus(\zeta_a^{(h)})\exp(-\mu_i^{(h)})>0$.
    $E^{\mathrm{ZBL}}$ is $C^\infty$ on $\Omega$.
    A finite composition and sum of $C^3$ scalars is $C^3$.
\end{proof}

\begin{theorem}[Conservativeness]\label{thm:cons}
    The forces $\bmF_i=-\partial E/\partial\bR_i$ and the stress $\bm{\Pi}_{\alpha\beta}=\partial E/\partial\epsilon_{\alpha\beta}\big|_{\epsilon=0}$ are the exact gradient and stress of the same scalar potential $E$ defined in \eqref{eq:Etot}.
\end{theorem}
\begin{proof}
    By Theorem~\ref{thm:smooth}, $E$ is a $C^3$ scalar function of $R$ and of the cell.
    The conclusion is the definition of gradient and stress of a scalar potential.
\end{proof}

\paragraph{Frozen zone.}
Say that a pair $(j,k)$ is \emph{frozen} if $r_{jk}\leq\rin$.

\begin{theorem}[Strict freezing of the learned branch]\label{thm:freeze}
    Suppose configuration $(Z,R)\in\Omega$ contains a frozen pair $(j,k)$.
    Then for every variation $\delta R$ that preserves frozen-ness (i.e.\ $r_{jk}(R+\delta R)\leq\rin$),
    \begin{equation}
        \frac{\partial E^{\mathrm{NN}}}{\partial\bR_j}(Z,R)\;=\;\frac{\partial E^{\mathrm{NN}}}{\partial\bR_k}(Z,R)\;=\;0,
    \end{equation}
    and consequently the forces on $j$ and $k$ coincide with the pure ZBL forces:
    \begin{equation}
        \bmF_j=-\partial E^{\mathrm{ZBL}}/\partial\bR_j,\qquad
        \bmF_k=-\partial E^{\mathrm{ZBL}}/\partial\bR_k.
    \end{equation}
\end{theorem}
\begin{proof}
    By Lemma~\ref{lem:eta} applied to $j$ (with $k\in\Nb_{\mathrm{out}}(j)$ and $r_{jk}\leq\rin$), $\eta_j=0$; symmetrically $\eta_k=0$.
    Every neighbour message reaching any node, in any aggregation summed over a source index $s$, carries a multiplicative factor $\eta_s$ — by \eqref{eq:gie}, \eqref{eq:agg-plain}, \eqref{eq:attw}, and (when active) \eqref{eq:film}.
    By induction on the block index, no node feature at any depth contains a contribution sourced from $j$ or $k$.
    At the source side itself, the clamped distance map fixes $\rtilde(r_{jk})=\rin$ and the clamped direction $\brh_{jk}$ enters only as part of an $\eta_j$-gated message, which vanishes.
    Therefore $E^{\mathrm{NN}}$ is constant under any variation that keeps the pair frozen.
    The force statement is the negative gradient of \eqref{eq:Etot}.
\end{proof}

% =============================================================
\section{Discussion: Local Frame, Global Frame, and Clebsch--Gordan}\label{sec:discuss-frames}

This section addresses a natural question raised by the absence of
equivariant--equivariant tensor products in DPA4 (\S\ref{sec:frame},
"trivialisation of Clebsch--Gordan"): \emph{since DPA4 only ever multiplies an
    equivariant feature by an $\SO(3)$-invariant scalar, do we still need the
    per-edge local frame, or could the entire architecture be re-expressed in the
    fixed global frame?} The answer is that the local frame is not needed for
avoiding Clebsch--Gordan, but it is needed for a more subtle purpose:
\emph{cross-degree, direction-dependent, edge-conditioned linear maps} on
equivariant features. This section makes the equivalence precise.

\subsection{Three CG-free regimes}\label{sec:three-regimes}

All three constructions below are $\SO(3)$-equivariant and use no
Clebsch--Gordan coefficient anywhere in their parameterisation; they differ
only in what class of operators they realise.

\paragraph{Regime (I): global frame, degree-preserving.}
The architecture is restricted to the global-frame operators of \S\ref{sec:ops}
(channel mixing $L^{\mathrm{ch}}$, degree-wise mixing $L^{\mathrm{deg}}$,
equivariant RMS norm, gated activation) plus invariant-weighted scatter
aggregation. The most general message $\bm{m}_{ij}^{\mathrm{(I)}}\in
    V_{\le\lmax}\otimes\R^{C}$ producible at edge $(i,j)$ within this regime is
\begin{equation}\label{eq:regI-msg}
    \bigl(\bm{m}_{ij}^{\mathrm{(I)}}\bigr)_{(l,m),c}
    \;=\;\sum_{c'}\bigl[\bm{W}^{(l)}(r_{ij};Z_i,Z_j)\bigr]_{c',c}\,\bh_{j,(l,m),c'},
\end{equation}
where $\bm{W}^{(l)}(\cdot)\in\R^{C\times C}$ is a per-degree, edge-conditioned channel-mixing matrix whose entries depend only on the $\SO(3)$-invariants $r_{ij},Z_i,Z_j$.
The aggregated update at atom $i$ is
\begin{equation}\label{eq:regI-agg}
    \bigl(\Delta\bh\bigr)_i \;=\; n_i\sum_{j:r_{ij}<\rcut} s(r_{ij})\,\eta_j\,\bm{m}_{ij}^{\mathrm{(I)}}.
\end{equation}
Note the strict block-diagonality in degree:
\begin{equation}\label{eq:regI-blockdiag}
    (\bm{m}_{ij}^{\mathrm{(I)}})_{(l_3,m_3)} \;=\;0\quad\text{whenever}\quad l_3\neq l_1\text{ for the source slot }l_1.
\end{equation}
There is no way for the source's $\bh_j^{(l_1)}$ to contribute to the destination's $\bh_i^{(l_3)}$ slot when $l_1\ne l_3$, nor any way for the message to depend on the edge direction $\brh_{ij}$.
Direction information enters only through the GIE seed of \eqref{eq:gie} at the initial state $X^{(0)}$.

\paragraph{Regime (II): local frame, DPA4.}
The architecture of \S\ref{sec:so2conv}: the source feature is transported to
the per-edge frame by $\bD_{ij}$, modulated by the scalar radial
$\widetilde\rho_{ij,l}$, mixed by the $\SO(2)_{\bz}$-equivariant linear
operator $L^{\mathrm{SO2}}$, and lifted back by
$K\,(\bD_{ij}^{\le\mmax})^{\top}$. The most general message
$\bm{m}_{ij}^{\mathrm{(II)}}$ producible at edge $(i,j)$ from a \emph{single}
source feature, with no bilinear products of two equivariant tensors, is
\begin{equation}\label{eq:regII-msg}
    \bm{m}_{ij}^{\mathrm{(II)}}\;=\;K\,\bigl(\bD_{ij}^{\le\mmax}\bigr)^{\top}\,L^{\mathrm{SO2}}\!\bigl[\,\diag(\widetilde\brho_{ij})\,\bigl(\bD_{ij}^{\le\mmax}\,\bh_j\bigr)\,\bigr],
\end{equation}
where $\diag(\widetilde\brho_{ij})$ is the diagonal matrix with $\widetilde\rho_{ij,l,c}$ on the $(l,\cdot,c)$ entry.
This composition can be \emph{not} block-diagonal in degree: the central operator $L^{\mathrm{SO2}}$ couples different $l$'s within each fixed $|m|$-stratum (Schur for $\SO(2)_{\bz}$, equation \eqref{eq:hom-so2}).
    We may rewrite \eqref{eq:regII-msg} in the equivalent global-frame form
    \begin{equation}\label{eq:regII-global}
        \bigl(\bm{m}_{ij}^{\mathrm{(II)}}\bigr)_{(l_3,m_3),c}
        \;=\;\sum_{l_1,m_1,c'}\bigl[\bm{\Phi}_{ij}\bigr]_{(l_3,m_3),(l_1,m_1),c',c}\,\bh_{j,(l_1,m_1),c'},
    \end{equation}
    where $\bm{\Phi}_{ij}$ is the edge-conditioned kernel
    \begin{equation}\label{eq:regII-Phi}
        \bm{\Phi}_{ij}\;=\;K\,\bigl(\bD_{ij}^{\le\mmax}\bigr)^{\top}\,\bm{W}\,\diag(\widetilde\brho_{ij})\,\bD_{ij}^{\le\mmax},
    \end{equation}
    with $\bm{W}$ collecting the parameters of $L^{\mathrm{SO2}}$.
    By construction $\bm{\Phi}_{ij}$ transforms equivariantly under $\SO(3)$ (rotations of the configuration), and its dependence on the edge direction $\brh_{ij}$ enters only through the two rotation operators $\bD_{ij}$.
    Aggregation is then \eqref{eq:regI-agg} with $\bm{m}_{ij}^{\mathrm{(II)}}$ in place of $\bm{m}_{ij}^{\mathrm{(I)}}$.

    \paragraph{Regime (III): global frame with Clebsch--Gordan, MACE/NequIP-style.}
    The angular content of the edge is carried explicitly as an equivariant feature
$Y_{l_2}^{m_2}(\brh_{ij})\in V_{l_2}$ for each $l_2\in\{0,\dots,\lmax\}$, and
    the message is built from a CG-tensor product of $\bh_j$ with
$\bm{Y}(\brh_{ij})$:
    \begin{equation}\label{eq:regIII-msg}
        \bigl(\bm{m}_{ij}^{\mathrm{(III)}}\bigr)_{(l_3,m_3),c}
        \;=\;\sum_{l_1,l_2}\bm{R}^{(l_1,l_2,l_3)}_{c}(r_{ij};Z_i,Z_j)\;\sum_{m_1,m_2}\bigl\langle l_1\,m_1;\,l_2\,m_2\,\big|\,l_3\,m_3\bigr\rangle\,\bh_{j,(l_1,m_1),c}\,Y_{l_2}^{m_2}(\brh_{ij}),
    \end{equation}
    with selection rule $|l_1-l_2|\le l_3\le l_1+l_2$ and learnable per-channel radial coefficients $\bm{R}^{(l_1,l_2,l_3)}_c(\cdot)\in\R$.
    This is a genuine bilinear equivariant product of two equivariant objects, $\bh_j$ and $\bm{Y}(\brh_{ij})$, and the Clebsch--Gordan symbols are unavoidable for the SO(3)-equivariant parameterisation.
    DPA4 does \emph{not} use this regime.

    \subsection{From Regime II to Regime III in the local frame}\label{sec:II-to-III}

    The passage from Regime II to Regime III is not obvious from the global-frame
    formulas above. In Regime II we used an $\SO(2)$-equivariant mixer
$L^{\mathrm{SO2}}$ surrounded by $\bD_{ij}$ rotations and never wrote a
    Clebsch--Gordan coefficient anywhere; in Regime III the Clebsch--Gordan
    coefficient is the defining ingredient. Yet both yield the same operator class.
    The cleanest way to see the equivalence is to write Regime III in the local
    frame and watch the Clebsch--Gordan coefficients trivialise. We do that
    explicitly here.

    \paragraph{The local-frame value of the angular feature.}
    The Regime III message is built from the spherical harmonic
$Y_{l_2}^{m_2}(\brh_{ij})$, viewed as an element of $V_{l_2}$. In the per-edge
    local frame, $\brh_{ij}$ is mapped to the local $\bz$-axis, so the relevant
    value is $Y_{l_2}^{m_2}(\bz)$. In the real-basis convention,
    \begin{equation}\label{eq:Y-at-pole}
        Y_{l_2}^{m_2}(\bz)\;=\;\sqrt{\frac{2l_2+1}{4\pi}}\,\delta_{m_2,0},
    \end{equation}
    i.e.\ the spherical harmonic at the north pole is a $\delta$-function on the $m_2=0$ line: every $m_2\ne 0$ component vanishes.
    All of the angular information that was carried by the $(2l_2+1)$-vector $\{Y_{l_2}^{m_2}(\brh_{ij})\}_{m_2}$ in the global frame is, in the local frame, concentrated entirely on the $m_2=0$ line; the rest of the information has been transferred to the rotation operator $\bD_{ij}$ that performed the frame change.

    \paragraph{The CG product in the local frame.}
    Start with the Regime III message \eqref{eq:regIII-msg} and rotate both the
    source feature and the angular feature into the per-edge local frame. Let
$\bh^{\mathrm{loc}}_j\coloneqq\bD_{ij}\bh_j$ denote the source feature in the
    local frame. The Regime III message in the local frame reads
    \begin{align}\label{eq:regIII-loc-1}
        \bigl[\bm{m}_{ij}^{\mathrm{(III)}}\bigr]^{\mathrm{loc}}_{(l_3,m_3),c}
        \; & =\;\sum_{l_1,l_2}\bm{R}^{(l_1,l_2,l_3)}_c(r_{ij})\sum_{m_1,m_2}\bigl\langle l_1\,m_1;\,l_2\,m_2\,\big|\,l_3\,m_3\bigr\rangle\,\bh^{\mathrm{loc}}_{j,(l_1,m_1),c}\,Y_{l_2}^{m_2}(\bz)                         \\
        \; & =\;\sum_{l_1,l_2}\bm{R}^{(l_1,l_2,l_3)}_c(r_{ij})\sqrt{\frac{2l_2+1}{4\pi}}\sum_{m_1}\bigl\langle l_1\,m_1;\,l_2\,0\,\big|\,l_3\,m_3\bigr\rangle\,\bh^{\mathrm{loc}}_{j,(l_1,m_1),c},\label{eq:regIII-loc-2}
    \end{align}
    where in the second line we substituted \eqref{eq:Y-at-pole} and collapsed the sum over $m_2$ to its only non-zero term $m_2=0$.

    \paragraph{$m$-conservation collapses the inner sum.}
    The Clebsch--Gordan coefficient satisfies the selection rule $m_1+m_2=m_3$, so
    for $m_2=0$ the only contributing $m_1$ is $m_1=m_3$. The inner sum in
    \eqref{eq:regIII-loc-2} reduces to a single term, and we may rewrite the entire
    message as
    \begin{equation}\label{eq:regIII-loc-3}
        \boxed{\;\bigl[\bm{m}_{ij}^{\mathrm{(III)}}\bigr]^{\mathrm{loc}}_{(l_3,m),c}
            \;=\;\sum_{l_1}\,\bm{W}^{(l_1,l_3,|m|)}_{c}(r_{ij})\,\bh^{\mathrm{loc}}_{j,(l_1,m),c},\;}
    \end{equation}
    where the effective per-edge mixer is the Clebsch--Gordan-weighted radial combination
    \begin{equation}\label{eq:W-from-R}
        \bm{W}^{(l_1,l_3,|m|)}_{c}(r_{ij})\;\coloneqq\;\sum_{l_2=|l_1-l_3|}^{l_1+l_3}\bm{R}^{(l_1,l_2,l_3)}_c(r_{ij})\sqrt{\frac{2l_2+1}{4\pi}}\,\bigl\langle l_1\,m;\,l_2\,0\,\big|\,l_3\,m\bigr\rangle.
    \end{equation}
    The CG coefficient $\langle l_1\,m;l_2\,0|l_3\,m\rangle$ is a fixed numerical constant; it depends on $m$ but not on the configuration.
    The channel index $c$ is preserved across input and output: \eqref{eq:regIII-loc-3} is the natural per-channel form inherited from the per-channel radials of \eqref{eq:regIII-msg}, with no cross-channel mixing.

    \paragraph{Eq.~\eqref{eq:regIII-loc-3} in the DPA4 edge-convolution notation.} The locally-rotated source feature
$\bh^{\mathrm{loc}}_j=\bD_{ij}\bh_j$ used in \eqref{eq:regIII-loc-3} is
    identical to the locally-rotated equivariant feature
$\bx_{ij}=\bD^{\le\mmax}_{ij}\bh_j$ of Stage 2 \eqref{eq:xloc} (modulo the
    harmless restriction $\bD\to\bD^{\le\mmax}$ that drops $|m|>\mmax$ components
    when the truncated layout is used). Substituting
$\bh^{\mathrm{loc}}_{j,(l_1,m),c}=\bx_{ij,(l_1,m),c}$ rewrites the Regime III
    message in the same input form as the DPA4 edge convolution pipeline:
    \begin{equation}\label{eq:regIII-loc-x}
        \bigl[\bm{m}_{ij}^{\mathrm{(III)}}\bigr]^{\mathrm{loc}}_{(l_3,m),c}\;=\;\sum_{l_1}\bm{W}^{(l_1,l_3,|m|)}_c(r_{ij})\,\bx_{ij,(l_1,m),c}.
    \end{equation}
    This is the cross-$l$-mixing generalisation of the per-degree radial modulation \eqref{eq:radmod}: where \eqref{eq:radmod} multiplies $\bx_{ij,(l,m),c}$ by a per-degree scalar $\widetilde\rho_{ij,l,c}$ that is diagonal in $l$, \eqref{eq:regIII-loc-x} replaces that diagonal scalar by an $|m|$-stratified matrix $\bm{W}^{(l_1,l_3,|m|)}_c$ that couples different degrees within the same $|m|$-stratum.
    The per-degree modulation \eqref{eq:radmod} is then the special case
    \begin{equation}\label{eq:radmod-as-W}
        \bm{W}^{(l_1,l_3,|m|)}_c(r_{ij})\;=\;\widetilde\rho_{ij,l_1,c}\,\delta_{l_1,l_3},
    \end{equation}
    i.e.\ a degree-diagonal mixer.
    Stages 3 and 5 of \S\ref{sec:so2conv} combined --- per-degree radial modulation \eqref{eq:radmod} followed by the $\SO(2)$-equivariant linear mixer $L^{\mathrm{SO2}}$ acting per $|m|$-stratum --- produce exactly the general $\bm{W}^{(l_1,l_3,|m|)}_c$ form of \eqref{eq:regIII-loc-x}, with $L^{\mathrm{SO2}}$ supplying the off-diagonal-in-$l$ couplings that \eqref{eq:radmod} cannot.
    The DPA4 edge convolution is therefore the local-frame version of the Regime III message construction, with the Clebsch--Gordan coefficients of \eqref{eq:W-from-R} absorbed into the learnable entries of $L^{\mathrm{SO2}}$ and the per-degree radial modulation $\widetilde\rho$.

    \paragraph{What \eqref{eq:regIII-loc-3} says.} Two features fall out of the local-frame rewriting:
    \begin{enumerate}[leftmargin=2em]
    \item \emph{The message is diagonal in $m$.}
    Source $m_1=m_3$: there is no cross-$m$ mixing at all.
    This is because the angular feature in the local frame has all its content at $m_2=0$, and $m$-conservation forces $m_1=m_3$.
    \item \emph{Within each fixed $|m|$-stratum, $l_1$ and $l_3$ are coupled arbitrarily.}
    The mixer $\bm{W}^{(l_1,l_3,|m|)}(r)$ is, for each fixed $(l_1,l_3,|m|)$, a numerical scalar (per channel pair) whose value depends on $r_{ij}$ through the radial coefficients $\bm{R}^{(l_1,l_2,l_3)}_c$ and on $m$ through the fixed CG coefficients.
    \end{enumerate}
    These are exactly the two features of Regime II's $\SO(2)$-equivariant operator $L^{\mathrm{SO2}}$ of \eqref{eq:so2-mn}: $|m|$-stratified, cross-$l$ mixing within each $|m|$-stratum, no cross-$|m|$ mixing.
    \emph{Regime III in the local frame is Regime II.}

    \paragraph{The two parameterisations side by side.}
    At the per-channel level, the local-frame mixer entry
$\bm{W}^{(l_1,l_3,|m|)}_c(r)$ admits two equivalent presentations, differing
    only in how the $r$- and $m$-dependence is decomposed.

    \medskip
    \noindent
    \emph{Regime II (per-channel).}
$\bm{W}^{(l_1,l_3,|m|)}_c(r)$ is an unconstrained learnable radial function, one per quadruple $(l_1,l_3,|m|,c)$.
    The $m$-dependence is independent across $|m|$-strata, and the total free-parameter count is one radial per $(l_1,l_3,|m|,c)$.

    \medskip
    \noindent
    \emph{Regime III (in the local frame).}
    The same per-edge mixer entry decomposes as
    \begin{equation}\label{eq:W-vs-R}
        \bm{W}^{(l_1,l_3,|m|)}_c(r)\;=\;\sum_{l_2=|l_1-l_3|}^{l_1+l_3}\bm{R}^{(l_1,l_2,l_3)}_c(r)\,K^{(l_1,l_2,l_3,m)},\qquad
        K^{(l_1,l_2,l_3,m)}\;\coloneqq\;\sqrt{\tfrac{2l_2+1}{4\pi}}\,\langle l_1\,m;\,l_2\,0\,|\,l_3\,m\rangle,
    \end{equation}
    with one learnable radial $\bm{R}^{(l_1,l_2,l_3)}_c(r)$ per CG triple $(l_1,l_2,l_3,c)$ and fixed numerical constants $K^{(l_1,l_2,l_3,m)}$ tying the $|m|$-strata together.

    \medskip
    The two presentations realise the same family of $|m|$-stratified per-channel mixers; they differ only in how the $m$-dependence is parameterised.
    Regime II treats different $|m|$-strata as independent radials; Regime III ties them together through a $(2\min(l_1,l_3)+1)$-dimensional basis of CG coefficients in $l_2$.

    \paragraph{A remark on channel mixing.}
    The discussion above is on the per-channel level (input channel $c$ maps to
    output channel $c$). The full DPA4 $L^{\mathrm{SO2}}$ operator additionally
    mixes channels via a $C\times C$ matrix per $(l_1,l_3,|m|)$-stratum, so Regime
    II as implemented is in fact the channel-mixing generalisation
$\bm{W}^{(l_1,l_3,|m|)}_{c,c'}(r)$ of the per-channel form above. The
    channel-mixing story is \emph{orthogonal} to the angular structure: it can be
    added to either Regime II or Regime III without affecting the $m$-conservation,
$|m|$-stratification, or CG-trivialisation conclusions of this subsection. Most
    MACE/NequIP-style Regime III networks add it through a final per-degree channel
    mixer $L^{\mathrm{ch}}$ following the CG product, exactly as in
    \eqref{eq:effn-std} of \S\ref{sec:effn}.

    \paragraph{The trivialisation of Clebsch--Gordan.}
    The reason no CG coefficient appears in DPA4 is now visible: when the angular
    feature $Y_{l_2}^{m_2}(\brh_{ij})$ is evaluated in the local frame, only its
$m_2=0$ component is non-zero (eq.~\eqref{eq:Y-at-pole}). The CG sum over
$(m_1,m_2)$ in the global formula collapses to a single $(m_1,0)$ term
    constrained by $m_1=m_3$, and the remaining $l_2$-sum becomes a CG-weighted sum
    of radial functions \eqref{eq:W-from-R} that can be absorbed into a single
    learnable mixer entry $\bm{W}^{(l_1,l_3,|m|)}(r)$ — exactly the entry of
$L^{\mathrm{SO2}}$ that DPA4 already provides. \emph{There is nothing to
    compute: the CG coefficients $\langle l_1\,m;l_2\,0|l_3\,m\rangle$ are fixed
    numerical constants that get folded into the learnable
    radial-times-channel-mixer parameters.} DPA4 chooses the direct
    parameterisation by $\bm{W}^{(l_1,l_3,|m|)}(r)$ entries; it never names the
    implicit $\bm{R}^{(l_1,l_2,l_3)}(r)$ that a Regime III architecture would name.

    \paragraph{What about $Y_{l_2}^{m_2\ne 0}(\brh_{ij})$, did we lose information?}
    The answer is no: the $m_2\ne 0$ components of $Y_{l_2}$ in the global frame are exactly the entries of the rotation operator $\bD_{ij}^{(l_2)}$ at the column corresponding to $\bz$.
    Rotating to the local frame moves this angular information from the value of $Y_{l_2}$ at $\brh_{ij}$ to the entries of $\bD_{ij}$ that bracket $L^{\mathrm{SO2}}$ in Regime II.
    The total content is preserved; only the placement of the angular information in the parameterisation changes.

    % =============================================================
    \subsection{Expressiveness ladder}\label{sec:expr-ladder}

    The three regimes form a strict inclusion of operator classes at the per-edge
    level:
    \begin{equation}\label{eq:ladder}
        \underbrace{\bigl\{\bm{m}_{ij}^{\mathrm{(I)}}\bigr\}}_{\text{degree-preserving}}
        \;\subsetneq\;
        \underbrace{\bigl\{\bm{m}_{ij}^{\mathrm{(II)}}\bigr\}}_{\text{cross-}l,\ \text{single-source}}
        \;\subseteq\;
        \underbrace{\bigl\{\bm{m}_{ij}^{\mathrm{(III)}}\bigr\}}_{\text{full SO(3)-equivariant bilinear}}.
    \end{equation}

    \begin{lemma}[Regime I is a proper subset of Regime II]
    Every message of the form \eqref{eq:regI-msg} is obtainable as a special case of \eqref{eq:regII-global} by choosing the cross-$l$ blocks of $\bm{W}$ in $L^{\mathrm{SO2}}$ to vanish.
    Conversely, since $\bm{\Phi}_{ij}$ in \eqref{eq:regII-Phi} generically has nonzero off-diagonal-in-$l$ entries (whenever $L^{\mathrm{SO2}}$'s cross-$l$ mixers are nonzero), Regime II achieves messages that violate \eqref{eq:regI-blockdiag}, so the inclusion is strict.
    \end{lemma}

    \paragraph{Why $l$-mixing is impossible in Regime I.}
    The block-diagonality \eqref{eq:regI-blockdiag} of Regime I messages is not an
    arbitrary stylistic choice; it is forced by Schur's lemma. For each fixed edge
    invariants $(r_{ij},Z_i,Z_j)$ the operator
    \[
    \bm{T}_{ij}:\bh_j\;\longmapsto\;\bm{m}_{ij}^{(\mathrm{I})}
\]
is an $\SO(3)$-equivariant linear map $V_{\le\lmax}\otimes\R^C\to
    V_{\le\lmax}\otimes\R^C$, because in Regime I the kernel depends only on
rotation-invariant data. Applying the real Schur classification of
\eqref{eq:hom-so3}, the space of such equivariant linear maps is
\[
    \mathrm{Hom}_{\SO(3)}\bigl(V_{\le\lmax}\otimes\R^C,V_{\le\lmax}\otimes\R^C\bigr)
    \;=\;\bigoplus_{l=0}^{\lmax}\R^{C\times C},
\]
i.e.\ one $C\times C$ matrix per degree, and the zero map between different
degrees: cross-$l$ entries are forbidden by Schur and not by an architectural
design choice. Equivalently in component form,
\begin{equation}\label{eq:schur-blockdiag}
    \bigl[\bm{T}_{ij}\bigr]_{(l_3,m_3),(l_1,m_1)}\;=\;\delta_{l_3,l_1}\,\delta_{m_3,m_1}\,\bm{W}^{(l_1)}(r_{ij};Z_i,Z_j),
\end{equation}
recovering exactly the form \eqref{eq:regI-msg}.
The only way to obtain a non-zero cross-$l$ matrix element while preserving $\SO(3)$-equivariance is to let the kernel depend on the edge direction $\brh_{ij}$ — and to do so in a way that transforms covariantly with the configuration under $\SO(3)$ rotations.
Regime II achieves this through the rotation operators $\bD_{ij},\bD_{ij}^{\top}$ that bracket $L^{\mathrm{SO2}}$ in \eqref{eq:regII-Phi}; Regime III achieves the same through the equivariant edge feature $\{Y_{l_2}^{m_2}(\brh_{ij})\}$ and an explicit Clebsch--Gordan contraction.
In both cases the kernel is no longer a fixed equivariant linear map but a \emph{family} of equivariant linear maps parameterised by $\brh_{ij}\in\Sph$, and Schur's no-cross-$l$ theorem of \eqref{eq:hom-so3} does not apply to such direction-dependent families.

\begin{lemma}[Regime II as a sub-parameterisation of Regime III]
    For any edge $(i,j)$, the kernel $\bm{\Phi}_{ij}$ of \eqref{eq:regII-Phi} admits a Clebsch--Gordan expansion of the form
    \begin{equation}\label{eq:Phi-CG-expand}
        \bigl[\bm{\Phi}_{ij}\bigr]_{(l_3,m_3),(l_1,m_1)}
        \;=\;\sum_{l_2=|l_3-l_1|}^{l_3+l_1}\widetilde R^{(l_1,l_2,l_3)}(r_{ij})\;\sum_{m_2}\bigl\langle l_1\,m_1;\,l_2\,m_2\,\big|\,l_3\,m_3\bigr\rangle\,Y_{l_2}^{m_2}(\brh_{ij}),
    \end{equation}
    where $\widetilde R^{(l_1,l_2,l_3)}(r_{ij})$ is an effective radial coefficient determined by $L^{\mathrm{SO2}}$ and $\widetilde\brho_{ij}$ via a finite linear combination of $\bm{W}$ entries.
    Hence Regime II messages are a subset of Regime III messages, obtained by tying together the $\widetilde R^{(l_1,l_2,l_3)}$ across different $l_2$ that share the same residual $|m|$-stratum.
\end{lemma}

The proof of the second lemma is a routine consequence of the matrix-element
factorisation $\bD_{ij}^{(l)}(R)=\sum Y_l^m(\hat{R\,\bz})\cdots$ together with
the Wigner-Eckart theorem; we omit it.

\subsection{Why DPA4 stops at Regime II}\label{sec:why-regII}

The architectural choice between (II) and (III) is not about expressiveness in
the limit but about \emph{parameterisation cost} of an equivariant cross-$l$
kernel:

\begin{itemize}[leftmargin=2em]
    \item In Regime III, every admissible triple $(l_1,l_2,l_3)$ with $|l_1-l_2|\le
              l_3\le l_1+l_2$ contributes one radial coefficient
          $\bm{R}^{(l_1,l_2,l_3)}_c(r)$ per channel, and the bilinear contraction over
          $(m_1,m_2)$ via the 3-$j$ symbols $\langle l_1 m_1;l_2 m_2|l_3 m_3\rangle$ has
          nominal cost $\mathcal{O}(\lmax^6)$ per edge.
    \item In Regime II, the same kernel is parameterised by the entries of
          $L^{\mathrm{SO2}}$, of size $\mathcal{O}(\lmax\,\mmax\,C^2)$, with the angular
          dependence carried implicitly by the rotation operators
          $\bD_{ij},\bD_{ij}^{\top}$ that act on the source feature. The per-edge cost is
          $\mathcal{O}(\lmax\,\mmax\,H)$ for the linear part and $\mathcal{O}(D\cdot
              D_m)$ for the rotation operators.
    \item Truncation by $\mmax\le\lmax$ in Regime II discards exactly those
          Clebsch--Gordan triples that would couple $|m|>\mmax$ channels in Regime III,
          recovering them only via the per-degree rescale $\rho_l$ of \eqref{eq:rescale}.
\end{itemize}

For DPA4's target operating regime ($\lmax$ small, $\mmax$ small, $C$ moderate,
$E$ large), the difference between the two parameterisations is dramatic:
Regime II's cost is essentially the same as a degree-preserving Regime-I
operator (linear in $\lmax$ and $\mmax$), while Regime III's cost grows as the
sixth power of $\lmax$. This is the algorithmic motivation for the local-frame
design.

\subsection{What the local frame is, mathematically}\label{sec:lf-summary}

To summarise the abstract content of \S\ref{sec:so2conv} in light of
\eqref{eq:regII-Phi}, the per-edge message $\bm{m}_{ij}^{\mathrm{(II)}}$ in
DPA4 is:
\begin{equation}\label{eq:lf-form}
    \boxed{\;
        \bm{m}_{ij}^{\mathrm{(II)}} \;=\; \bm{\Phi}_{ij}\,\bh_j,\qquad
        \bm{\Phi}_{ij}\;\coloneqq\;K\,\bigl(\bD_{ij}^{\le\mmax}\bigr)^{\top}\,\bm{W}\,\diag(\widetilde\brho_{ij})\,\bD_{ij}^{\le\mmax}
        \;}
\end{equation}
— a per-edge $\SO(3)$-equivariant linear map between two equivariant features, applied to a single source.
No object on the right-hand side is the tensor product of two equivariant features; both rotations $\bD_{ij},\bD_{ij}^{\top}$ act on the same single feature $\bh_j$, and all scalars ($\widetilde\brho_{ij}$, $K$) are invariants of the edge.
The Clebsch--Gordan coefficients that would explicitly expand $\bm{\Phi}_{ij}$ in the basis $\{Y_l^m(\brh_{ij})\}$, as in \eqref{eq:Phi-CG-expand}, are implicit in the matrix elements of $\bD_{ij}$ and never need to be computed in floating-point arithmetic.

The local frame is therefore best understood not as a different ingredient from
Regime III but as a different \emph{parameterisation} of the same
SO(3)-equivariant edge-conditioned linear-map space, organised to make the
angular dependence cheap and the cross-$l$ mixing rich. If one were to drop the
local frame entirely (Regime I), one would lose the cross-$l$ mixing but retain
everything else, including the equivariance, smoothness, and freezing
properties of \S\ref{sec:thms}.

% =============================================================
\section*{Notation Index}

\begin{tabular}{ll}
    $N$                                                                                                & number of atoms                                             \\
    $Z_i\in\{1,\dots,n_t\}$                                                                            & atomic species of atom $i$                                  \\
    $\bR_i\in\R^3$                                                                                     & position of atom $i$                                        \\
    $\br_{ij}=\bR_j-\bR_i$, $r_{ij}=\|\br_{ij}\|$, $\brh_{ij}=\br_{ij}/r_{ij}$                         & edge vector, distance, direction                            \\
    $\lmax,\mmax\in\Z_{\geq 0}$                                                                        & maximum degree, maximum truncation order $(\mmax\leq\lmax)$ \\
    $D=(\lmax+1)^2$, $D_m=\sum_l(2\min(l,\mmax)+1)$                                                    & full and truncated dimensions                               \\
    $\iota(l,m)=l^2+l+m$                                                                               & linear packing of $(l,m)$ in $V_{\le\lmax}$                 \\
    $V_l$, $D^l$, $V_{\le\lmax}$, $\bD(R)$                                                             & irreducible blocks; full Wigner-D matrix                    \\
    $\bq_{ij}\in S^3$, $R_{ij}=R(\bq_{ij})$, $\bD_{ij}=\bD(R_{ij})$                                    & edge quaternion and rotation                                \\
    $\bP_{\mmax}$, $K$ with entries $\rho_l$                                                           & truncation projector and per-degree rescale                 \\
    $s(r)=s_5(r)$, $\sigma_r(r)=s_7(r)$                                                                & cutoff envelopes                                            \\
    $\rin<\rout$, $\rtilde(r)$, $w(r)$                                                                 & bridging window and switches                                \\
    $\eta_j=\prod_{i\in\Nb_{\mathrm{out}}(j)} w(r_{ji})$                                               & source-freeze gate                                          \\
    $d_i=\sum_j s(r_{ij})^2$, $n_i=(d_i+\varepsilon)^{-1/2}$                                           & smooth degree and its inverse square root                   \\
    $\bphi(r)\in\R^{n_r}$                                                                              & trainable-frequency radial basis                            \\
    $\Te_{\mathrm{edge}}(i,j)=\Te^{(1)}_{Z_i,:}+\Te^{(2)}_{Z_j,:}\in\R^{C}$                            & per-edge species-pair embedding                             \\
    $\brho_{ij}\in\R^{(\lmax+1)\times C}$                                                              & per-edge per-degree radial features (\eqref{eq:rho-edge})   \\
    $C, F, C_f, H=FC_f$                                                                                & channel, focus, per-focus widths                            \\
    $\bh_i\in V_{\le\lmax}\otimes\R^C$                                                                 & node feature at atom $i$                                    \\
    $L^{\mathrm{ch}}$, $L^{\mathrm{deg}}$, $L^{\mathrm{SO2}}$, $N_{\gamma,\beta}$, $\Gamma_{\psi,\bG}$ & equivariant operators \S\ref{sec:ops}                       \\
    $\mathcal{C}_\theta$, $\mathcal{F}_\theta$, $\mathcal{B}_\theta$                                   & edge convolution, EFFN, interaction block                   \\
    $E^{\mathrm{NN}}$, $E^{\mathrm{ZBL}}$, $E$                                                         & learned, analytical, and total energies
\end{tabular}

\end{document}
